% concatenated from main_arxiv.tex, src/formatting.tex, src/macros.tex, sections/00-abstract.tex, sections/01-intro.tex, sections/02-lower-bounds.tex, sections/03-sdp-discretization.tex, sections/04-computations.tex, sections/05-disk-annulus.tex, sections/06-conclusion.tex, sections/07-fourier-expansion.tex
\documentclass[a4paper, 10pt]{amsart}
\usepackage[margin=3.75cm]{geometry}
\usepackage{amsmath,amssymb,amsthm}
\usepackage{esint}
\usepackage{amsaddr}
\usepackage{graphicx}
\usepackage{booktabs}
\usepackage[sort&compress,numbers]{natbib}
\usepackage[colorlinks=true]{hyperref}
\usepackage[capitalize,nameinlink]{cleveref}

%% Formatting and macros in other files
%% Equation numbering
\numberwithin{equation}{section}

%% Captions
\usepackage[font=footnotesize,labelfont=bf,width=\textwidth]{caption}

%% Font size in references
\def\bibfont{\small}

%% Link colors
\hypersetup{
    linkcolor = black, 
    filecolor = blue,
    urlcolor  = blue,
    citecolor = blue,
    }
\urlstyle{same}

%% Lists
\usepackage[shortlabels]{enumitem}

%% Cleveref equation names
\crefname{equation}{}{}

%% Theorems
\theoremstyle{plain}
\newtheorem{theorem}{Theorem}[section]
\newtheorem{proposition}[theorem]{Proposition}
\newtheorem{lemma}[theorem]{Lemma}
\newtheorem{corollary}[theorem]{Corollary}
%
\theoremstyle{definition}
\newtheorem{definition}{Definition}[section]
\newtheorem{assumption}{Assumption}[section]
%
\theoremstyle{remark}
\newtheorem{remark}{Remark}[section]

%% Fix for cleveref names?
%% Did not need this in the past...PROBLEM WITH NEW LATEX VERSION
\AddToHook{env/lemma/begin}{\crefalias{theorem}{lemma}}
\AddToHook{env/proposition/begin}{\crefalias{theorem}{proposition}}
\AddToHook{env/corollary/begin}{\crefalias{theorem}{corollary}}

%% Cleveref names
% \crefname{assumption}{assumption}{assumptions}
% \Crefname{assumption}{Assumption}{Assumptions}
% \crefname{lemma}{lemma}{lemmas}
% \Crefname{lemma}{Lemma}{Lemmas}
% \crefname{proposition}{proposition}{propositions}
% \Crefname{proposition}{Proposition}{Propositions}
% \crefname{figure}{figure}{figures}
% \Crefname{figure}{Figure}{Figures}

% Bold subsection numbers
% define what you want for the various levels
\makeatletter
\def\@seccntformat#1{%
  \protect\textup{%
    \protect\@secnumfont
    \expandafter\protect\csname format#1\endcsname % <--- added
    \csname the#1\endcsname
    \protect\@secnumpunct
  }%
}
\newcommand{\formatsubsection}{\bfseries}
% Spaces
\newcommand{\cal}[1]{\mathcal{#1}}
\newcommand{\R}{\mathbb{R}}
\newcommand{\N}{\mathbb{N}}
\newcommand{\Z}{\mathbb{Z}}
\newcommand{\Pe}{\mathrm{Pe}}
\newcommand{\Hpp}{\cal{H}^{++}}
\newcommand{\Hpm}{\cal{H}^{+-}}
\newcommand{\Hmp}{\cal{H}^{-+}}
\newcommand{\Hmm}{\cal{H}^{--}}
\newcommand{\BMO}{\rm{BMO}(\Omega)}
\newcommand{\Hardy}{\mathcal{H}^1{(\Omega)}}
\newcommand{\U}{\mathcal{U}}
\newcommand{\Uh}{\mathcal{U}_h}

\newcommand{\Hmo}{H^{-1}}
\newcommand{\Hoz}{H_0^1(\Omega)}
\newcommand{\Ltwo}{L^2(\Omega)}

\newcommand{\nxi}{n}
\newcommand{\ntheta}{n}
\newcommand{\npsi}{n}

\DeclareMathOperator{\imag}{Im}
\renewcommand{\Im}{\imag}

% Optimization problem objects
\newcommand{\QF}{\mathcal{Q}}
\newcommand{\cost}[2]{\avg{f#1 - \frac14 \abs{\nabla #1}^2} - #2}

% Math operators
\DeclareMathOperator{\divergence}{div}

% Other math
\renewcommand{\vec}[1]{\mathbf{#1}}
\newcommand{\weakto}{\rightharpoonup}
\newcommand{\abs}[1]{\left\vert #1 \right\vert}
\newcommand{\dx}{\mathrm{d} \vec{x}}
\newcommand{\dy}{\mathrm{d} \vec{y}}
\newcommand{\dz}{\mathrm{d} \vec{z}}
\newcommand{\dr}{\mathrm{d} r}
\newcommand{\dtheta}{\mathrm{d} \theta}
\renewcommand{\div}{\divergence}%{\nabla \cdot }
\newcommand{\contembbed}{\hookrightarrow}
\newcommand{\invlap}{\Delta^{-1}}
\newcommand{\ndof}{N_{\mathrm{dof}}}
\newcommand{\avg}[1]{\int_\Omega #1 \,\dx}
\newcommand{\norm}[1]{\left\| #1 \right\|}
\newcommand{\e}{{\rm e}}

% Optimal value and lower bound
\newcommand{\ocpval}{\mathcal{J}}
\newcommand{\efficiency}{\mathcal{E}}
\newcommand{\lb}{\mathcal{L}}

% Convenient text stuff
\newcommand{\Peclet}{P\'eclet}
\newcommand{\comment}[1]{{\color{red}[#1]}}


%% Article info
\title[Optimizing bounds for energy-constrained cooling in 2D]{Optimizing bounds for energy-constrained optimal cooling problems in two dimensions}
\author{Pedro Bl\"oss Braga and Giovanni Fantuzzi}
\address{
    Department of Mathematics\\
    Friedrich--Alexander Universit\"at Erlangen--N\"urnberg
}
\email{\href{mailto:pedro.bloess@fau.de}{pedro.bloess@fau.de}}
\email{\href{mailto:giovanni.fantuzzi@fau.de}{giovanni.fantuzzi@fau.de}}


\thanks{GF is supported by the DFG project MONET (grant no. 568735456). PB is supported by the DFG project \emph{Renormalized Solutions in der Optimalsteuerung von Evolutionsgleichungen} (grant no. 551486487).}

\date{\today}
%%%%%%%%%%%%%%%%%%%%%%%%%%%%%%%%%%%%%%%%%%%%%%%%%%%%%%%%%%%%%%%%%%%%%%%%%%%%%%%%
%%%%%%%%%%%%%%%%%%%%%%%%%%%%%%%%%%%%%%%%%%%%%%%%%%%%%%%%%%%%%%%%%%%%%%%%%%%%%%%%
\begin{document}
\begin{abstract}
We study optimal control problems for incompressible fluids in two-dimensional domains with a cold boundary and internal heat sources and sinks. Given a kinetic energy budget, measured by the square of a nondimensional \Peclet\ number $\Pe$, the goal is to maximize a cooling efficiency inversely proportional to the mean square gradient of the fluid's temperature.
Using Lagrange duality, we formulate a well-posed dual problem whose solution yields an upper bound on the maximum cooling efficiency $\efficiency(\Pe)$. We then numerically approximate the dual problem using a convergent hierarchy of semidefinite programs obtained via discretization.
We illustrate this computational approach on optimal cooling problems in a square and in an annulus, explaining also how problem symmetries can be exploited to reduce computational complexity. Finally, we construct admissible points for the dual problem to obtain new analytical upper bounds on the optimal cooling efficiency $\efficiency(\Pe)$. Specifically, we prove that $\efficiency(\Pe) \lesssim \Pe^{2}$ for arbitrary domains and heat distributions, and that $\efficiency(\Pe) \lesssim \Pe^{2}/ \ln^2\Pe$ for cooling flows in disks and annuli with heat source/sink distributions with a positive azimuthal average. These results generalize and improve known efficiency bounds for energy-constrained cooling flows in a disk.
\end{abstract}

%% Title and contents
\maketitle

%%%%%%%%%%%%%%%%%%%%%%%%%%%%%%%%%%%%%%%%%%%%%%%%%%%%%%%%%%%%%%%%%%%%%%%%%%%%%%%%
\section{Introduction}
Optimal cooling problems, in which the flow of a fluid must be controlled to efficiently cool domains with internal heat sources, 
are fundamental in fluid mechanics. They are central to the design of heat exchangers and also provide a paradigm for studying transport processes.
Locally optimal flows can usually be computed with algorithms for PDE-constrained optimization (see, e.g., \cite{Hassanzadeh2014,Alben2017a,Alben2017b,Motoki2018,Souza2020}), and provide lower bounds on the optimal cooling efficiency. Upper bounds, in contrast, are substantially more difficult to compute because they are statements about the efficiency of all possible flows.
Remarkably, the difficulty persists even in a highly simplified setting where the admissible fluid flows are steady, heat is a ‘passive’ scalar with no effect on the fluid's motion, and the only constraints on the flow are incompressibility and a bound on some norm of the fluid's velocity. In this work, we focus on two-dimensional incompressible flows with prescribed $L^2$ norm, which physically correspond to flows with a prescribed kinetic energy.

\subsection{Problem description}
Consider an incompressible fluid flowing inside a bounded, possibly multiply connected domain $\Omega \subset \R^2$ with Lipschitz boundary $\partial\Omega$. The flow is described by a time-independent and weakly divergence-free velocity field $\vec{u} \in L^2(\Omega; \R^2)$ satisfying the no-penetration condition $\vec{u}\cdot \hat{\vec{n}}=0$ on $\partial\Omega$, where $\hat{\vec{n}}$ is the outward-pointing unit vector normal to the boundary.
The domain is subject to a distribution $f \in H^{-1}(\Omega)$ of heat sources and sinks, while its boundary is held at a constant temperature of zero. The temperature $T$ of the fluid then solves the steady advection-diffusion equation
\begin{equation}\label{e:ade}
    \begin{cases}
        - \Delta T + \vec{u}\cdot \nabla T =  f &\text{on }\Omega,\\
        T=0 &\text{on }\partial\Omega.
    \end{cases}
\end{equation}

We wish to understand how well the domain can be cooled by flows that do not exceed a prescribed kinetic energy budget, measured by a nondimensional parameter $\Pe$ called the \Peclet\ number. We consider a flow efficient at cooling if the temperature $T$ is nearly constant on $\Omega$, meaning that $\smash{\avg{\abs{\nabla T}^2}}$ is small. Then, the optimal flow solves the minimization problem
\begin{equation}\label{e:ocp}\tag{OCP}
    \ocpval(\Pe):=
    \min_{\substack{
        \vec{u}\in L^2(\Omega;\R^2) \\ 
        \div\vec{u}=0 \text{ on }\Omega\\ 
        \vec{u} \cdot \hat{\vec{n}} = 0  \text{ on }\partial\Omega\\
        \norm{\vec{u}}_{L^2} \leq \,\Pe
    }}\;
    \avg{\abs{\nabla T}^2}
\end{equation}
This is a nonconvex optimal control problem for the advection-diffusion equation \cref{e:ade}, where the nonconvexity is due entirely to the bilinear term $\vec{u}\cdot\nabla T$ in this equation. We associate to the optimal control cost $\ocpval(\Pe)$ an optimal cooling efficiency via the formula
\begin{equation}\label{e:efficiency}
    \efficiency(\Pe) := \frac{\|f\|_{\Hmo}^2}{ \ocpval(\Pe)},
\end{equation}
where the norm $\|f\|_{\Hmo}$ is included to have $\efficiency(0) = 1$. In other words, we measure the optimal cooling efficiency relative to the no-flow state. 

The definitions immediately imply that $\efficiency(\Pe) \geq 1$. Better lower bounds can be found by constructing particular flows, including locally optimal ones computed with PDE-constrained optimization algorithms. Our aim is to obtain complementary upper bounds by deriving positive lower bounds on $\ocpval(\Pe)$.

\subsection{Known bounds}
When $\Omega$ is a disk and the heating is uniform ($f=1$), numerical optimization produces convection rolls achieving $\smash{\avg{\abs{\nabla T}^2}\sim \Pe^{-1}}$~\cite{Marcotte2018}. Inspired by these computations, analytical constructions for a broad class of sufficiently smooth heating distributions $f$ provide the lower bound $\efficiency(\Pe) \gtrsim \Pe$ at sufficiently large \Peclet\ numbers, with a prefactor depending only on the heating distribution~\cite[Proposition~5.1]{Tobasco2022}.
This estimate is in fact true for a time-dependent generalization of our optimal cooling problem, which we do not consider for simplicity. 

Lower bounds on optimal cooling efficiencies have also been derived for variations of problem \cref{e:ocp}, such as: optimal cooling in a disk under the `enstrophy' constraint $\smash{\norm{\vec{u}}_{H^1}\leq\Pe}$ \cite{Tobasco2022}; optimal cooling of uniformly heated flows between an insulated boundary and an isothermal boundary, with constraints on the $L^p$ and $W^{1,p}$ norms of $\vec{u}$ \cite{Iyer2022}; optimal cooling in insulated domains with balanced heat sources and sinks \cite{Shaw2007,Song2023}; and optimal cooling between differentially heated plates \cite{Hassanzadeh2014,TobascoDoering2017,DoeringTobasco2019,Kumar2024}. In all cases, lower bounds on the cooling efficiency are derived from particular flows that resemble those obtained with numerical optimization \cite{Motoki2018,Souza2020,Marcotte2018}.

Finding good upper bounds on $\efficiency(\Pe)$ has proven more difficult. The problem has received attention within a broader study of heat transport and mixing of passive scalars (see, e.g., \cite{Doering1996,Hassanzadeh2014,Motoki2018,TobascoDoering2017,DoeringTobasco2019,Souza2020,Kumar2024,Chanillo2025,Song2023,Shaw2007,Thiffeault2004,DoeringThiffeault2006,Thiffeault2012}) and upper bounds on cooling efficiencies have been derived for many of the settings described above \cite{Tobasco2022,Shaw2007,Song2023,Kumar2024}. However, they often overestimate the corresponding lower bounds. In particular, the upper bound $\efficiency(\Pe) \lesssim \Pe^2$ for the energy-constrained problem \cref{e:ocp} in a disk~\cite[Proposition~2.6]{Tobasco2022} does not match the lower bound $\efficiency(\Pe) \gtrsim \Pe$. It is then natural to wonder if the proofs of these upper bounds can be improved and, in particular, whether numerical computations can be used to guide such improvements.

\subsection{Contributions}
We combine Lagrange duality with semidefinite programming to optimize lower bounds on the optimal control cost $\ocpval(\Pe)$, which translate into upper bounds on the cooling efficiency $\efficiency(\Pe)$ via \cref{e:efficiency}. Our key insight is to recognize that existing proofs of lower bounds for $\ocpval(\Pe)$ implicitly construct feasible points for the Lagrange dual of problem \cref{e:ocp}, which is an infinite-dimensional semidefinite program (SDP). We explicitly derive this dual problem, prove its well-posedness (\cref{prop:ldp-well-posed}), and show that it can be approximated with arbitrary accuracy by finite-dimensional SDPs (\cref{thm:convergence}). These are computationally tractable for fixed $\Pe$, so the best bound provable via Lagrange duality can be computed numerically.

We also use the dual of \cref{e:ocp} to derive new analytical upper bounds on $\efficiency(\Pe)$. First, we give an elementary proof that
\begin{equation}\label{e:an-bound-intro}
    \efficiency(\Pe) \leq 1 + C \Pe^2
    \qquad \forall\, \Pe \geq 0,
\end{equation}
where $C>0$ depends only on the fluid's domain $\Omega$ (\cref{th:analytical-lb}). This result applies to arbitrary heat distributions $f$ and domains $\Omega$, generalizing a similar bound proved in \cite{Tobasco2022} when $\Omega$ is a disk. 
We then fix $\Omega$ to be a disk or an annulus and prove that, when $f$ has a strictly positive average on any circle around the domain's center,
\begin{equation}\label{e:intro-log-lb}
    \efficiency(\Pe) \leq C \, \frac{\Pe^2}{\ln^2 \Pe}
    \qquad \forall\, \Pe > 0
\end{equation}
with a constant $C$ depending on the domain and on $f$ (\cref{th:refined-annulus,th:refined-disk}). This logarithmically-improved bound holds in particular if the heating is uniform ($f=1$). Its proof is strongly inspired by the numerical results produced by our semidefinite programming computations for a uniformly heated annulus.

\subsection{Relation to prior work}
Lagrange duality is a classical approach for computing bounds on optimization problems. While it is often left implicit, it underpins most existing bounds on the cooling efficiency or other heat transport properties of incompressible flows. In particular, duality arguments similar to ours have been applied to optimal cooling problems in insulated domains with balanced heating and cooling \cite{Song2023}. What distinguishes our work is that we explicitly formulate the Lagrange dual of problem \cref{e:ocp} in a form suitable for a numerical implementation based on SDPs.

The connection between optimal cooling problems and SDPs is not surprising: discretizations of problem \cref{e:ocp} are nonconvex quadratic programs, and the Lagrange dual of a quadratic program is an SDP~\cite{Shor1987,Nesterov2000}. This fact has been exploited to compute performance bounds for quadratic-program discretizations of optimal design problems in structural mechanics and photonics~\cite{Tyburec2019,Angeris2019,Angeris2021,Gertler2025,Dalklint2025}. Our approach is similar, but we dualize before discretizing. Exchanging these operations has a twofold advantage: we can prove that SDP solutions converge to the best bound provable with Lagrange duality, and we can derive rigorous bounds analytically. 

Finally, our approach is similar to SDP-based implementations of the \emph{background method}, a variational framework for bounding the time-averaged dissipation in incompressible flows (see \cite{Doering1994,Constantin1995,Doering1996}, as well as the review \cite{Fantuzzi2022} and the many references therein). This is because the variational problem at the heart of the background method is the Lagrange dual of a quadratic program relaxation of the optimization problem for the time-averaged dissipation over flows satisfying the incompressible Navier--Stokes equations (see~\cite{Kerswell1998,Kerswell1999} for details).

\subsection{Outline}
The plan of the paper is as follows. In \cref{s:dual-problem}, we derive and analyze the Lagrange dual of problem \cref{e:ocp}. 
In particular, we construct an admissible point for the dual problem to obtain an analytical lower bound on $\ocpval(\Pe)$, which combined with \cref{e:efficiency} leads to the efficiency bound in \cref{e:an-bound-intro}. In \cref{s:discretization}, we introduce SDP discretizations of the Lagrange dual problem and prove that they converge with increasing resolution for all discretization schemes with strong approximation properties in the Sobolev space $H^1(\Omega)$. This convergence is not unexpected, but the proof requires specialized arguments due to the non-standard, `infinite-dimensional SDP' structure of the Lagrange dual problem. 
We demonstrate our numerical strategy in \cref{s:numerics} on energy-constrained cooling problems in a square and in an annulus, outlining also how problem symmetries can be exploited to reduce computational complexity. 
In \cref{s:improved-bounds}, inspired by our computations, we provide a refined analysis of optimal cooling in disks and annuli, proving the efficiency bound in \cref{e:intro-log-lb}. \Cref{s:conclusion} offers final remarks and perspectives for future work.
\section{The Lagrange dual problem}\label{s:dual-problem}

In this section, we formulate and study the Lagrange dual of the optimal cooling problem \cref{e:ocp}. To obtain this dual problem in a convenient form for both analysis and computations, in \cref{ss:streamfunction} we reformulate \cref{e:ocp} in terms of a scalar streamfunction. We derive the dual of the reformulated problem in \cref{ss:dualization}, study its optimizers in \cref{ss:well-posedness}, and use it to deduce the efficiency bound \cref{e:an-bound-intro} in \cref{ss:an-lower-bound}.


%%%%%%%%%%%%%%%%%%%%%%%%%%%%%%%%%%%%%%%%%%%%%%%%%%%%%%%%%%%%%%%%%%%%%%%%%%%%%%%%
\subsection{Streamfunction formulation}\label{ss:streamfunction}
Let $\Gamma_0$ denote the exterior boundary of the domain $\Omega$ and, if this is multiply connected, let $\Gamma_1,\ldots,\Gamma_p$ denote all other connected parts of the boundary (see \cref{f:domain-with-holes} for an illustration). Let $\cal{H}$ be the subspace of $H^1(\Omega)$ contaning functions that are constant on each piece of the boundary, that is,
\begin{equation*}
    \cal{H} := \left\{ \psi \in H^1(\Omega):\;\exists c_0,\ldots,c_p \in \R \text{ s.t. } \psi|_{\Gamma_i} = c_i ,\; 0\leq i \leq p \right\}.
\end{equation*}
Every weakly divergence-free vector field $\vec{u}\in L^2(\Omega;\R^2)$ satisfying the no-penetration boundary condition $\vec{u}\cdot \hat{\vec{n}}=0$ can be expressed as
\begin{equation}\label{e:streamfunction-def}
    \mathbf{u} = \Pe\, \nabla^\perp \psi = \Pe\left( \partial_y \psi,\, -\partial_x \psi\right)
\end{equation}
for a streamfunction $\psi \in \cal{H}$ determined up to an additive constant (see, e.g., \cite[\S3.1]{GiraultRaviart1979}). We fix this constant such that $\psi|_{\Gamma_0} = 0$, so $\cal{H}=H^1_0(\Omega)$ if $\Omega$ is simply connected. The factor of $\Pe$ in \cref{e:streamfunction-def} is included for convenience.

\begin{figure}[t]
    \centering
    \includegraphics[scale=0.9]{img/domain-sketch.png}
    \caption{Illustration of a multiply connected domain $\Omega$ (shaded gray region) with exterior boundary $\Gamma_0$ (thick gray line) and interior boundaries $\Gamma_1$ (thick green line) and $\Gamma_2$ (thick red line).}
    \label{f:domain-with-holes}
\end{figure}

With this re-parametrization, the optimal cooling problem \cref{e:ocp} becomes
\begin{equation}\label{e:ocp-psi}
    \ocpval(\Pe)=
    \min_{\substack{
        \psi \in \cal{H} \\ 
        \norm{\nabla \psi}_{L^2} \leq 1
    }}\;
    \avg{\abs{\nabla T}^2},
\end{equation}
where $T$ is the unique weak solution of the advection-diffusion equation \cref{e:ade} with velocity field $\vec{u}=\Pe \nabla^\perp \psi$. Specifically, $T$ is the unique function in $H^1_0(\Omega)$ satisfying
\begin{equation}\label{e:ade-weak}
    \avg{ \nabla \xi \cdot \nabla T + \Pe\,\xi \,\nabla^\perp \psi \cdot \nabla T - f\xi}= 0 \quad \forall \xi \in H^1_0(\Omega).
\end{equation}
The existence and uniqueness of such a weak solution follows from standard arguments based on the Lax--Milgram theorem; the only difficulty is to show that the product $\xi \,\nabla^\perp \psi \cdot \nabla T$ is integrable despite the low regularity of the factors, but this is true because $\nabla^\perp \psi \cdot \nabla T = \partial_y\psi \partial_x T - \partial_x \psi \partial_y T$ is a Jacobian determinant. In particular, as discussed in  \cite[Remark~2.2]{Tobasco2022} and \cite[\S3.1]{Song2023}, analysis in \cite{Coifman1990} implies that
\begin{equation}\label{e:clms}
    \avg{ \xi \,\nabla^\perp \psi \cdot \nabla T } \leq c(\Omega) \norm{\nabla \xi}_{L^2} \norm{\nabla \psi}_{L^2} \norm{\nabla T}_{L^2}
\end{equation}
for a positive constant $c(\Omega)$ depending only on the domain $\Omega$.
This estimate will play a key role also in our numerical analysis of \cref{ss:convergence}. 


%%%%%%%%%%%%%%%%%%%%%%%%%%%%%%%%%%%%%%%%%%%%%%%%%%%%%%%%%%%%%%%%%%%%%%%%%%%%%%%%
\subsection{The dual problem}\label{ss:dualization}
We now dualize the minimization problem in \cref{e:ocp-psi} to obtain a convex maximization problem whose optimal value, denoted by $\lb(\Pe)$, is a lower bound on $\ocpval(\Pe)$.  To ease the notation, we associate to every pair $(a,\xi) \in \R \times H^1_0(\Omega)$ a homogeneous quadratic form $\mathcal{Q}(a,\xi):\cal{H} \times H^1_0(\Omega) \to \R$ defined as
\begin{equation}\label{e:qf}
    \mathcal{Q}\left(a,\xi\right)[\psi,\theta] := \avg{ a|\nabla \psi|^2 + \abs{\nabla\theta}^2 - \Pe\,\xi\, \nabla^\perp \psi \cdot \nabla \theta}.
\end{equation}
Note that the map $(a,\xi) \mapsto \mathcal{Q}(a,\xi)$ is affine and that the integral is well defined thanks to estimate \cref{e:clms}. We write $\mathcal{Q}\left(a,\xi\right)\succeq 0$ to indicate that $\mathcal{Q}\left(a,\xi\right)$ is a positive semidefinite quadratic form, meaning that $\mathcal{Q}\left(a,\xi\right)[\psi,\theta] \geq 0$ for all $\psi \in \cal{H}$ and $\theta \in H^1_0(\Omega)$.

\begin{proposition}\label{th:duality}
    For every $\Pe\geq 0$, 
    the optimal value $\ocpval(\Pe)$ of problem \cref{e:ocp} satisfies $\ocpval(\Pe) \geq \lb(\Pe) \geq 0$, where
    \begin{equation}\label{e:ldp}\tag{LDP}
        \lb(\Pe) := \max_{\substack{
            a\in \R \\ 
            \xi \in H^1_0(\Omega) \\ 
            \QF\left(a,\xi\right) \succeq  0 
        }}\;
        \avg{f\xi - \frac14 \abs{\nabla \xi}^2} - a.
    \end{equation}
\end{proposition}

\begin{remark}
    For $\Pe > 0$, we expect the inequality $\ocpval(\Pe) \geq \lb(\Pe)$ to be strict in general because the optimal cooling problem \cref{e:ocp} is not convex. For $\Pe=0$, both \cref{e:ocp} and \cref{e:ldp} are trivially solvable and attain the same value (see also \cref{rem:Pe=0} below).
\end{remark}

\begin{remark}\label{rem:symmetrization-1}
    The constraint $\QF\left(a,\xi\right) \succeq  0$ requires $a\geq 0$ and we can in fact restrict the optimization in \cref{e:ldp} to $a>0$. We can then set $a=\alpha^2$ with $\alpha > 0$ and perform the change of variables $\psi \mapsto \psi / \sqrt{\alpha}$, $\theta \mapsto \theta \sqrt{\alpha}$ to conclude that $\QF\left(a,\xi\right)$ is positive semidefinite if and only if so is the `balanced' quadratic form
    \begin{equation}\label{e:qf-symmetrized}
    \mathcal{Q}'\left(\alpha,\xi\right)[\psi,\theta] := \avg{ \alpha|\nabla \psi|^2 + \alpha\abs{\nabla\theta}^2 - \Pe\,\xi\, \nabla^\perp \psi \cdot \nabla \theta}.
    \end{equation}
    Consequently,
    \begin{equation*}
        \lb(\Pe) = \max_{\substack{
            \alpha \in \R \\ 
            \xi \in H^1_0(\Omega) \\ 
            \QF'\left(\alpha,\xi\right) \succeq  0 
        }}\;
        \avg{f\xi - \frac14 \abs{\nabla \xi}^2} - \alpha^2.
    \end{equation*}
    We will use this reformulation in \cref{s:improved-bounds} to prove \cref{e:intro-log-lb}. We also remark that the feasible set of this balanced problem formulation is a cone in $\R\times\Hoz$, but exploring the computational implications of this fact is beyond the scope of our present work.
\end{remark}

\begin{remark}\label{rem:symmetrization-2}
    When $\Omega$ is simply connected, $\cal{H}=\Hoz$ and the balanced quadratic form $\QF'(a,\xi)$ satisfies $\mathcal{Q}'\left(a,\xi\right)[\psi,\theta] = \mathcal{Q}'\left(a,\xi\right)[\theta,-\psi]$. As demonstrated in \cref{ss:example-square}, this identity can afford computational savings if combined with other problem symmetries.
\end{remark}

\begin{proof}[Proof of \cref{th:duality}]
    The temperature $T$ in \cref{e:ocp} is the unique function in $H^1_0(\Omega)$ satisfying \cref{e:ade-weak}. Viewing the test function $\xi$ in that equation as a Lagrange multiplier for the advection-diffusion equation \cref{e:ade}, and introducing a Lagrange multiplier $a\geq 0$ for the constraint $\norm{\nabla\psi}_{L^2}\leq 1$, we can write
    \begin{align*}
        \ocpval(\Pe)
        &= 
        \inf_{\substack{
        \psi \in \cal{H} \\ 
        T \in H^1_0(\Omega)
        }}\;
        \sup_{\substack{a \in \R \\ \xi \in H^1_0(\Omega) \\ a\geq 0}}
        \avg{\abs{\nabla T}^2 + a |\nabla \psi|^2 + (\nabla\xi - \Pe\,\xi\,\nabla^\perp\psi )\cdot \nabla T + f\xi} - a
        \\
        &\geq 
        \sup_{\substack{a \in \R \\ \xi \in H^1_0(\Omega) \\ a\geq 0}}
        \inf_{\substack{
        \psi \in \cal{H} \\ 
        T \in H^1_0(\Omega)
        }}\;
        \avg{\abs{\nabla T}^2 + a |\nabla \psi|^2 + (\nabla\xi - \Pe\,\xi\,\nabla^\perp\psi ) + f\xi} - a.
    \end{align*}
    The inequality on the second line is due to switching minimization and maximization.

    Next, we perform the inner minimization over $T$ for fixed $\psi$, $a$ and $\xi$. The optimal temperature $T^*$ solves the Euler--Lagrange equation
    $-2\Delta T^* - \Delta \xi + \Pe \, \nabla \xi \cdot \nabla^\perp \psi = 0$
    with boundary condition $T^* = 0$. Writing $\invlap$ for the inverse Laplacian with homogeneous Dirichlet boundary conditions, we obtain
    \begin{equation}\label{e:Topt}
        T^* = - \frac12\xi + \frac{\Pe}{2}\ \invlap(\nabla \xi \cdot \nabla^\perp \psi).
    \end{equation}
    To compute the value of the minimum, we multiply the Euler--Lagrange equation by $T^*$ and integrate by parts over $\Omega$ to find that
    \begin{equation}\label{e:el-1}
    \avg{ (\nabla\xi - \Pe\,\xi\,\nabla^\perp\psi )\cdot \nabla T^*} = -2 \avg{\abs{\nabla T^*}^2}.
    \end{equation}
    We then observe that
    \begin{equation}\label{e:el-2}
        \avg{\nabla \xi \cdot \nabla \invlap(\nabla \xi \cdot \nabla^\perp \psi)}
        = -\avg{ \xi \nabla \xi \cdot \nabla^\perp \psi}
        = 0,
    \end{equation}
    where the both equalities follows from integration by parts using the divergence-free nature of $\nabla^\perp \psi$ and the homogeneous Dirichlet boundary conditions of $\xi$. Using \cref{e:Topt}, \cref{e:el-1} and \cref{e:el-2}, we obtain
    \begin{align*}
        \ocpval(\Pe)
        \geq 
        \sup_{\substack{a \in \R \\ \xi \in H^1_0(\Omega) \\ a\geq 0}}
        \inf_{\substack{ %% begin constraints
        \psi \in \cal{H}
        }}\; %% end constraints
        \avg{
            a |\nabla \psi|^2 
            - \frac14 \abs{\nabla \xi}^2 
            - \frac{\Pe^2}{4} \abs{\nabla \invlap(\nabla \xi \cdot \nabla^\perp \psi)}^2 
            + f\xi
        } 
        - a.
    \end{align*}

    We now remove the term involving the inverse Laplacian using the identity 
    \begin{equation*}
        -\avg{ \frac{\Pe^2}{4} \abs{\nabla \invlap(\nabla \xi \cdot \nabla^\perp \psi)}^2 } = \min_{\theta \in H^1_0(\Omega)} \avg{ \abs{\nabla\theta}^2 - \Pe\,\xi\,\nabla^\perp \psi \cdot \nabla\theta },
    \end{equation*}
    arriving at the lower bound
    \begin{align*}
        \ocpval(\Pe)
        \geq 
        \sup_{\substack{a \in \R \\ \xi \in H^1_0(\Omega) \\ a\geq 0}}
        \inf_{\substack{ %% begin constraints
        \psi \in \cal{H} \\ 
        \theta \in H^1_0(\Omega)
        }}\; %% end constraints
        \avg{a\abs{\nabla\psi}^2 + \abs{\nabla\theta}^2 - \Pe\,\xi\,\nabla^\perp\psi \cdot \nabla\theta - \frac14 \abs{\nabla \xi}^2 + f\xi} - a.
    \end{align*}
    The minimax problem on the right-hand side of this inequality is equivalent to the maximization problem \cref{e:ldp} because the inner minimum is zero if the quadratic form $\QF(a,\xi)$ defined in \cref{e:qf} is positive semidefinite, and is negative infinity otherwise. Therefore, the outer maximization selects pairs $(a,\xi)$ satisfying $\QF(a,\xi)\succeq 0$. Since this constraint automatically implies that $a\geq 0$, the latter condition need not be enforced in \cref{e:ldp} and we conclude that $\ocpval(\Pe) \geq \lb(\Pe)$. Finally, $\lb(\Pe)\geq 0$ because the choices $a=0$ and $\xi=0$ are feasible for \cref{e:ldp} and have a zero objective value.
\end{proof}


%%%%%%%%%%%%%%%%%%%%%%%%%%%%%%%%%%%%%%%%%%%%%%%%%%%%%%%%%%%%%%%%%%%%%%%%%%%%%%%%
\subsection{Optimizers and strictly feasible near-optimizers}\label{ss:well-posedness}
In stating problem \cref{e:ldp}, we tacitly assumed that the maximum is attained. We prove this fact next, establishing also that the maximizer is unique.

\begin{proposition}\label{prop:ldp-well-posed}
    The maximization problem \cref{e:ldp} has a unique optimal solution.
\end{proposition}

\begin{proof}
    Let $\{(a_n,\xi_n)\}_{n\in\N}$ be a maximizing sequence for \cref{e:ldp}. The semidefiniteness of $\QF(a_n,\xi_n)$ implies that $a_n\geq 0$. We may also assume that $\smash{\cost{\xi_n}{a_n}} \geq 0$ because $\lb(\Pe)\geq 0$. Then, using the H\"older, Poincar\'e and Young inequalities, we can find $C>0$ such that
    \begin{align}\label{e:boundedness}
        a_n + \frac14\norm{\nabla \xi_n}_{L^2}^2 
        \leq C \norm{f}_{\Hmo} \norm{\nabla \xi_n}_{L^2} 
        \leq 2C^2\norm{f}_{\Hmo}^2 + \frac18\norm{\nabla \xi_n}_{L^2}^2
    \end{align}
    for every $n$. This inequality implies that the sequence $\{(a_n,\xi_n)\}_{n\in\N}$ is uniformly bounded, so we may replace it with a subsequence converging to a limit $(a^*,\xi^*)$ weakly in $\R \times H^1_0(\Omega)$. 
    
    We claim that this limit is an optimizer for \cref{e:ldp}. To see that it is feasible, recall from \cref{e:qf} that the map $(a,\xi)\mapsto \QF(a,\xi)$ is affine, so for every $\psi\in\cal{H}$ and every $\theta \in H^1_0(\Omega)$ we have
    \begin{equation*}
        \QF(a^*,\xi^*)[\psi,\theta] = \lim_{n\to\infty}  \underbrace{\QF(a_n,\xi_n)[\psi,\theta]}_{\geq 0} \geq 0
    \end{equation*}
    by the definition of weak convergence.
    For optimality, instead, observe that the function $(a,\xi)\mapsto \smash{\cost{\xi}{a}}$ is the sum of a linear function and a negative multiple of the squared $H^1_0$ norm, so it is weakly upper semicontinuous in $\R \times H^1_0(\Omega)$. Then,
    \begin{equation*}
        \lb(\Pe) = \lim_{n\to\infty} \cost{\xi_n}{a_n} \leq \cost{\xi^*}{a^*} \leq \lb(\Pe),
    \end{equation*}
    from which we conclude that $(a^*,\xi^*)$ is optimal for \cref{e:ldp}.

    Finally, we show that $(a^*,\xi^*)$ is the only optimal solution. The strict concavity of the cost function with respect to $\xi$ implies that if $(a',\xi')$ is another optimizer for \cref{e:ldp}, then $\xi'=\xi^*$. But then, we must also have $a'=a^*$ because the optimality of these constants for \cref{e:ldp} implies that
    \begin{align*}
        a' 
        = 
        \sup_{\psi,\theta}
        \frac{\avg{\Pe\,\xi'\,\nabla^\perp\psi\cdot\nabla\theta - \abs{\nabla\theta}^2}}{\avg{\abs{\nabla \psi}^2}}
        =
        \sup_{\psi,\theta}
        \frac{\avg{\Pe\,\xi^*\,\nabla^\perp\psi\cdot\nabla\theta - \abs{\nabla\theta}^2}}{\avg{\abs{\nabla \psi}^2}}
        = a^*,
    \end{align*}
    where the suprema are taken over $\psi \in \cal{H}$ and $\theta \in \Hoz$. 
\end{proof}


%%%%%%%%%%%%%%%%%%%%%%%%%%%%%%%%%%%%%%%%%%%%%%%%%%%%%%%%%%%%%%%%%%%%%%%%%%%%%%%%
Finally, with a view towards the numerical analysis of \cref{e:ldp} in \cref{ss:convergence}, we show that problem \cref{e:ldp} admits nearly optimal solutions for which the quadratic form $\QF(a,\xi)$ is positive definite.

\begin{proposition}\label{prop:near-optimizers}
    Fix $\Pe\geq 0$. For every $\varepsilon \in [0,1]$, there exist $a \in \R$ and $\xi \in H^1_0(\Omega)$ such that $\smash{\cost{\xi}{a} \geq \lb(\Pe) - \varepsilon}$ and such that
    \begin{gather*}
        \QF(a, \xi)[\psi,\theta] \geq  \frac{\varepsilon}{1+\lb(\Pe)}\avg{\abs{\nabla\psi}^2 + \abs{\nabla\theta}^2}
        \quad\forall(\psi,\theta) \in \cal{H} \times \Hoz.
    \end{gather*}
\end{proposition}
\begin{proof}
    Let $(a^*,\xi^*)$ be the optimizer for \cref{e:ldp}. For every $\delta \in [0,1]$, set $\xi = (1-\delta) \xi^*$ and $a = (1-\delta) a^* + \delta$. Straightforward calculations show that
    \begin{equation*}
        \cost{\xi}{a} = \lb(\Pe) - \delta \left(1 + \lb(\Pe) \right)
    \end{equation*}
    and
    \begin{align*}
        \QF(a_\varepsilon,\xi_\varepsilon)[\psi,\theta] 
        &= (1-\delta)\QF(a^*,\xi^*)[\psi,\theta] 
        + \delta \avg{\abs{\nabla\psi}^2 + \abs{\nabla\theta}^2}.
    \end{align*}
    Since $\QF(a^*,\xi^*)$ is positive semidefinite, the claim follows upon taking $\smash{\delta = \frac{\varepsilon}{1+\lb(\Pe)}}$, which is an admissible choice because $\lb(\Pe)\geq 0$.
\end{proof}


%%%%%%%%%%%%%%%%%%%%%%%%%%%%%%%%%%%%%%%%%%%%%%%%%%%%
\subsection{An explicit optimal cooling efficiency bound}\label{ss:an-lower-bound}
We conclude this section by using the dual problem \cref{e:ldp} to prove the optimal cooling efficiency bound in \cref{e:an-bound-intro}. The proof relies on the identities 
\begin{equation*}
    \norm{f}_{\Hmo} 
    := \sup_{\varphi \in \Hoz} \frac{\abs{\avg{f \varphi}}}{\avg{\abs{\nabla \varphi}^2}} 
    = \left( \avg{ \abs{ \nabla \invlap f }^2 } \right)^\frac12,
\end{equation*}
where $\invlap$ is the inverse Laplacian with homogeneous Dirichlet boundary conditions. 

\begin{theorem}\label{th:analytical-lb}
    Let $c=c(\Omega)$ be the constant in estimate \cref{e:clms}. For every $\Pe\geq 0$,
    \begin{equation}\label{e:an-bound-theorem}
        \ocpval(\Pe)\geq \lb(\Pe)\geq \frac{\|f\|_{\Hmo}^2}{ 1 + c^2 \Pe^2}.
    \end{equation}
    Consequently, $\efficiency(\Pe) \leq 1 + c^2 \Pe^2$.
\end{theorem}

\begin{remark}
    This result extends to general domains $\Omega$ and arbitrary $\Pe$ the lower bound on $\ocpval(\Pe)$ obtained in \cite[Proposition~2.6]{Tobasco2022} when $\Omega$ is a disk and $\Pe$ is sufficiently large. The analysis in \cite{Tobasco2022}, however, applies also to time-dependent optimal cooling problems. We expect a version of \cref{th:analytical-lb} to hold for time-dependent problems, too.
\end{remark}

\begin{remark}\label{rem:Pe=0}
    The inequalities in \cref{th:analytical-lb} are equalities if $\Pe=0$, because in this case problem \cref{e:ocp} is trivially solved by $\vec{u}=0$, giving $T=-\invlap f$ and thus $\ocpval(0) = \|f\|_{\Hmo}^2$.
\end{remark}

\begin{proof}
    The bound on $\efficiency(\Pe)$ follows from \cref{e:efficiency,e:an-bound-theorem}.
    The first inequality in \cref{e:an-bound-theorem} is ensured by \cref{th:duality}. To prove the second one, let us fix
    $\smash{\xi = -2 \invlap f /(1 + c^2 \Pe^2)}$ and $\smash{a = c^2 \Pe^2 \norm{f}_{\Hmo}^2 / (1 + c^2 \Pe^2)^2}$.
    The identities 
    \begin{equation}\label{e:xi2f}
        c \,\Pe \norm{\nabla \xi}_{L^2} = \frac{2 c \,\Pe}{1 + c^2 \Pe^2} \norm{ f }_{\Hmo} = 2\sqrt{a}
    \end{equation}
    imply that
    $a \norm{\nabla \psi}_{L^2}^2 + \norm{\nabla \theta}_{L^2}^2 - c\, \Pe \norm{ \nabla \xi }_{L^2} \norm{\nabla \psi}_{L^2} \norm{\nabla \theta}_{L^2} \geq 0$   
    because the expression on the left-hand side is a perfect square. This inequality ensures that the quadratic form $\QF(a,\xi)$ in \cref{e:qf} is positive definite by estimate \cref{e:clms}, so the chosen $a$ and $\xi$ are admissible for problem \cref{e:ldp}.
    We then conclude that
    \begin{align*}
        \lb(\Pe) 
        &\geq \avg{ f\xi - \frac14 \abs{\nabla\xi}^2} - a \\
        &= -\frac{2}{ 1 + c^2 \Pe^2} \avg{ f\invlap f} - \frac{\norm{ f }_{\Hmo}^2}{(1 + c^2 \Pe^2)^2} 
        - \frac{ c^2 \Pe^2 }{\left(1 + c^2 \Pe^2 \right)^2} \norm{f}_{\Hmo}^2\\
        &= \frac{\|f\|_{\Hmo}^2}{ 1 + c^2 \Pe^2}.
    \end{align*}
   Here, we used the first identity in \cref{e:xi2f} and the definition of $a$ to obtain the second line, and the identity $\smash{\avg{f \invlap f} = - \norm{f}_{\Hmo}^2}$ to obtain the last line.
\end{proof}
\section{SDP discretizations}\label{s:discretization}

We now show that the optimal solution of problem \cref{e:ldp} can be approximated numerically by solving a hierarchy of increasingly large SDPs obtained via discretization.

\subsection{A general discretization scheme}
To approximate problem \cref{e:ldp} numerically, we discretize the space $\Hoz$ in which the function $\xi$ is optimized as well as the space $\cal{H}\times \Hoz$ on which the quadratic form $\QF(a,\xi)$ must be positive semidefinite. We do so using a generic conforming Galerkin discretization scheme, which in practice could be a pseudo-spectral or a finite element method.

Specifically, for each $n \in \N$, we introduce $n$-dimensional subspaces\footnote{We take these subspaces to have the same dimension for notational simplicity, but there are no obstacles to considering subspaces with different dimensions. In this case, the convergence results stated in \cref{ss:convergence} apply with $n = \min\left\{\dim \Xi_n,\, \dim \Theta_n,\, \dim \Psi_n \right\}$.}
\begin{align*}
    \Xi_n &\subset \Hoz, &
    \Theta_n &\subset \Hoz, &
    \Psi_n &\subset \cal{H}.
\end{align*}
Then, we define the discrete version of \cref{e:ldp} to be
\begin{equation}\label{e:ldp-discrete}
    \lb_n(\Pe) := \max_{\substack{a\in\R \\ \xi_n \in \Xi_n \\ \QF_n(a,\xi_n)\succeq 0}} \
    \avg{f\xi_n - \frac14 \abs{\nabla \xi_n}^2} - a,
\end{equation}
where $\QF_n$ is the restriction of the quadratic form $\QF$ to $\Psi_n \times \Theta_n$ and the constraint $\QF_n(a,\xi_n)\succeq 0$ simply means that $\QF(a,\xi_n)[\psi_n, \theta_n]\geq 0$ for all $\psi_n \in \Psi_n$ and $\theta_n \in \Theta_n$. The existence of a unique maximizer for \cref{e:ldp-discrete} follows exactly as in \cref{prop:ldp-well-posed}, and $\lb_n(\Pe)\geq 0$ because the choices $a=0$ and $\xi_n = 0$ are admissible for \cref{e:ldp-discrete}.

We stress that $\lb_n(\Pe)$ only approximates the optimal value $\lb(\Pe)$ of \cref{e:ldp}, but is neither an upper bound nor a lower bound. This is because we restrict the feasible set of \cref{e:ldp} by optimizing $\xi$ only over the subspace $\Xi_n$, but at the same time we relax the constraint $\QF(a,\xi)\succeq 0$ by imposing it only over $\Psi_n \times \Theta_n$. Lower bounds on $\lb(\Pe)$ could be computed using null Lagrangian translations or dual `occupation measure' relaxations~\cite{Chernyavsky2023,Valmorbida2016,Korda2022}, but we do not presently know if these lower bounds converge. In contrast, as we prove in \cref{ss:convergence}, the approximations given by \cref{e:ldp-discrete} converge under mild and reasonable assumptions on the subspaces $\Xi_n$, $\Theta_n$ and $\Psi_n$.


\subsection{Explicit SDP formulation}
The discrete dual problem \cref{e:ldp-discrete} is an SDP with a quadratic objective function. To make this evident, we introduce a basis $\{u_1,\ldots,u_n\}$ for $\Xi_n$, a basis $\{v_1,\ldots,v_n\}$ for $\Theta_n$, a basis $\{w_1,\ldots,w_n\}$ for $\Psi_n$, and expand
\begin{align*}
        \xi_n(\vec{x}) &= \sum_{i=1}^{\nxi} z_i u_i(\vec{x}), &
        \theta_n(\vec{x}) &= \sum_{i=1}^{\ntheta} t_i v_i(\vec{x}), &
        \psi_n(\vec{x}) &= \sum_{i=1}^{\npsi} p_i w_i(\vec{x}).
\end{align*}
We collect the expansion coefficients in $n$-dimensional column vectors $\vec{z}$, $\vec{t}$ and $\vec{p}$, define the $n$-dimensional column vector $\vec{f}=(\smash{\avg{f u_1}},\ldots,\smash{\avg{f u_n}})$, and consider the $n \times n$ matrices $A$, $B$, $C$ and $D_1,\ldots,D_n$ with entries
\begin{align*}
    A_{jk} &= \avg{\nabla u_j \cdot \nabla u_k}, &
    B_{jk} &= \avg{\nabla v_j \cdot \nabla v_k}, \\
    C_{jk} &= \avg{\nabla w_j \cdot \nabla w_k}, &
    (D_i)_{jk} &= \avg{u_i \nabla v_k \cdot \nabla^\perp w_j}.
\end{align*}
Using these quantities, we can write
${\avg{f\xi_n} = \vec{f}^\top \vec{z} }$, ${\avg{ \abs{\nabla \xi_n}^2 } = \vec{z}^\top A \vec{z}}$, and
\begin{equation*}
    \QF(a,\xi_n)[\psi_n, \theta_n] 
    =
    \begin{bmatrix}
        \vec{p} \\ \vec{t}
    \end{bmatrix}^\top
    \begin{bmatrix}
        a B &\ -\frac\Pe2 \sum_{i} z_i D_i \\
        -\frac\Pe2 \sum_{i} z_i D_i^\top & C
    \end{bmatrix}
    \begin{bmatrix}
        \vec{p} \\ \vec{t}
    \end{bmatrix}.
\end{equation*}
This quantity is nonnegative for all $\vec{p}$ and $\vec{t}$ (equivalently, for all $\psi_n\in\Psi_n$ and $\theta_n \in \Theta_n$) if and only if the $2\times 2$ block matrix on the right-hand side is positive semidefinite. Therefore, the `abstract' discrete problem \cref{e:ldp-discrete} is equivalent to
\begin{align}\label{e:sdp}
    \max_{\substack{a \in \R \\ \vec{z}\in\R^n}} \quad
    &\vec{f}^\top \vec{z} - \frac14 \vec{z}^\top A \vec{z} - a
    \\\nonumber \text{s.t}\quad &
     \begin{bmatrix}
        a B &\ -\frac\Pe2 \sum_{i} z_i D_i \\
        -\frac\Pe2 \sum_{i} z_i D_i^\top & C
    \end{bmatrix}
    \succeq 0.
\end{align}
This is a semidefinite program with a concave quadratic objective and can be solved by software packages such as \texttt{Clarabel}~\cite{Clarabel2024}. Linear conic program solvers such as \texttt{Mosek}~\cite{mosek} can be used after lifting the quadratic term $\vec{z}^\top A \vec{z}$ into a second-order cone constraint.


\subsection{Convergence analysis}\label{ss:convergence}
We now prove that the optimal value and optimizer of the discrete dual problem \cref{e:ldp-discrete} converge to the optimal value and optimizer of \cref{e:ldp} as the dimension of the subspaces $\Xi_n$, $\Theta_n$ and $\Psi_n$ increases. The only requirement, stated precisely in \cref{ass:density} below, is the existence of `recovery sequences' with elements in these subspaces that approximate functions in $\Hoz$ and $\cal{H}$. This approximation property is typical of pseudo-spectral discretization spaces or finite element space.

\begin{assumption}[Existence of recovery sequences]\label{ass:density}
The following conditions hold:
\begin{enumerate}[a), noitemsep, widest={a)}, leftmargin=*]
    \item $\forall \xi \in \Hoz$, $\exists \{\xi_n\}_{n\in\N}$ with $\xi_n \in \Xi_n$ such that $\lim_{n\to\infty}\|\xi_n - \xi\|_{\Hoz} = 0$.
    \item $\forall \theta \in \Hoz$, $\exists \{\theta_n\}_{n\in\N}$ with $\theta_n \in \Theta_n$ such that $\lim_{n\to \infty}\|\theta_n - \theta\|_{\Hoz} = 0$.
    \item $\forall  \psi \in \cal{H}$, $\exists \{\psi_n \}_{n \in \N}$ with $\psi_n \in \Psi_n$ such that $\lim_{n\to\infty}\|\psi_n - \psi\|_{H^1(\Omega)} = 0$.
\end{enumerate}
\end{assumption}

To simplify the presentation, we also assume that the integrals required to set up problem \cref{e:ldp-discrete} can be computed exactly. This assumption is mild, and quadrature errors that decrease sufficiently fast with $n$ could be accounted for if desired. With these assumptions in place, we are ready to state and prove the main result of this section.

\begin{theorem}\label{thm:convergence}
   Fix $\Pe\geq 0$. Let $\lb(\Pe)$ and $(a^*,\xi^*)$ be the optimal value and optimizer of problem \cref{e:ldp}. For each $n\in \N$, let $\lb_n(\Pe)$ and $(a^*_n,\xi^*_n)$ be the optimal value and optimizer of problem \cref{e:ldp-discrete}. If \cref{ass:density} holds, then 
   \begin{enumerate}[{\rm a)}, noitemsep, widest={a)}, leftmargin=*]
       \item\label{item:L-conv}  $\lb_n(\Pe) \to \lb(\Pe)$ as $n\to\infty$,
       \item\label{item:a-conv}  $a_n^* \to a^*$ as $n\to\infty$,
       \item\label{item:xi-conv}  $\xi_n^* \to \xi^*$ strongly in $\Hoz$ as $n\to\infty$.
   \end{enumerate}
\end{theorem}

\begin{proof}
The estimates in \cref{e:boundedness} guarantee that the sequences $\{a^*_n\}$ and $\{\xi^*_n\}$ are uniformly bounded in $\R$ and $\Hoz$, respectively, so there exist subsequences $\smash{\{a^*_{n_i}\}}$ and $\smash{\{\xi^*_{n_i}\}}$ such that $\smash{a^*_{n_i}\to a'}$ and $\smash{\xi^*_{n_i}} \weakto \xi'$ weakly in $\Hoz$ for some limit points $a'$ and $\xi'$. We claim, and will prove later, that
\begin{subequations}
    \begin{gather}
        \label{e:claim1}
        \lb(\Pe) \leq \liminf_{n\to\infty} \lb_n(\Pe) \\
        \label{e:claim2}
        \QF(a',\xi')\succeq 0.
    \end{gather}
\end{subequations}
The second claim means that the pair $(a',\xi')$ is feasible for \cref{e:ldp}. Then, since the function $(a,\xi)\mapsto \smash{\cost{\xi}{a}}$ is weakly upper semicontinuous, we can start from \cref{e:claim1} to estimate
\begin{align}\label{e:liminf-limsup-combo}
    \lb(\Pe)
    &\leq \liminf_{n\to\infty} \lb_n(\Pe)\\ \nonumber
    &\leq \liminf_{n_i\to\infty} \cost{\xi_{n_i}^*}{a_{n_i}^*}\\ \nonumber
    &\leq \limsup_{n_i\to\infty} \cost{\xi_{n_i}^*}{a_{n_i}^*}\\ \nonumber
    &\leq \cost{\xi'}{a'}
    \\ \nonumber
    &
    \leq \lb(\Pe).\phantom{\int_\Omega}
\end{align}
These inequalities must be equalities, so the weak subsequence limits $a'$ and $\xi'$ are optimizers of \cref{e:ldp}. Since this problem has unique optimizers $a^*$ and $\xi^*$, and since the arguments above apply to every convergent subsequence of discrete optimizers, we conclude that $a_n^* \to a^*$ and $\xi^*_n \weakto \xi^*$ weakly in $\Hoz$. We also conclude that the inequalities in \cref{e:liminf-limsup-combo} hold without passing to subsequences, so $\lb_n(\Pe)\to\lb(\Pe)$. 

To complete the proof, we now need to prove the claimed inequalities \cref{e:claim1} and \cref{e:claim2}, as well as the strong convergence $\xi^*_n \to \xi^*$ in $\Hoz$. 

To prove strong convergence, we observe that the weak convergence $\smash{\xi^*_n \weakto \xi^*}$ and the limit $\smash{\lb_n(\Pe)\to\lb(\Pe)}$ imply that $\smash{\avg{\abs{\nabla\xi^*_n}^2}\to \avg{\abs{\nabla\xi^*}^2}}$. Then, since $\smash{\Hoz}$ is a uniformly convex Banach space, $\{\xi_n^*\}$ converges strongly (see, e.g., \cite[Proposition~3.32]{Brezis2011}).

Next, we prove inequality \cref{e:claim1}. Fix an arbitrary $\varepsilon \in (0,1)$ and let $(a_\varepsilon,\xi_\varepsilon)$ be a strictly feasible near-optimizer for \cref{e:ldp} constructed using \cref{prop:near-optimizers}. By \cref{ass:density}, there exist $\xi_n \in \Xi_n$ such that $\xi_n \to \xi_\varepsilon$ strongly in $\Hoz$. We claim that the pair $(a_\varepsilon, \xi_n)$ is feasible for the discrete dual problem \cref{e:ldp-discrete} for all sufficiently large $n$. Indeed, for every pair $(\psi,\theta) \in \Psi_n \times \Theta_n \subset \cal{H}\times \Hoz$ we can use \cref{e:clms} to estimate
\begin{align*}
    \QF(a_\varepsilon, \xi_n)[\psi,\theta] 
    &= \QF(a_\varepsilon, \xi_\varepsilon)[\psi,\theta] + \Pe \avg{(\xi_\varepsilon-\xi_n) \nabla^\perp \psi \cdot \nabla\theta}
    \\
    &\geq \frac{\varepsilon}{1+\mathcal{L(\Pe)}} \avg{\abs{\nabla\psi}^2 + \abs{\nabla\theta}^2} - 
    c \Pe \norm{\nabla \xi_\varepsilon - \nabla \xi_n}_{L^2} \norm{\nabla \psi}_{L^2} \norm{\nabla\theta}_{L^2}
    \\
    &\geq \left(\frac{\varepsilon}{1+\mathcal{L(\Pe)}} - \frac{c\Pe}{2} \norm{\nabla \xi- \nabla \xi_n}_{L^2} \right)\avg{\abs{\nabla\psi}^2 + \abs{\nabla\theta}^2}
\end{align*}
and last quantity is positive for large $n$ because $\xi \to \xi_n$ in $\Hoz$. We then estimate
\begin{align*}
    \liminf_{n \to \infty}\lb_n(\Pe)
    &\geq \liminf_{n \to \infty} \cost{\xi_n}{a_\varepsilon} 
    \\
    &= \cost{\xi_\varepsilon}{a_\varepsilon}
    \geq \lb(\Pe) - \varepsilon
\end{align*}
and let $\varepsilon \to 0$ to obtain \cref{e:claim1}.

Finally, we establish \cref{e:claim2}. Fix $\psi \in \cal{H}$ and $\theta \in \Hoz$. We need to show that $\QF(a', \xi')[\psi, \theta] \geq 0$. To do this, we use \cref{ass:density} to construct `recovery sequences' $\{\psi_n\}$ and $\{\theta_n\}$ with $\psi_n\in \Psi_n$ and $\theta_n \in \Theta_n$ that converge strongly to $\psi$ and $\theta$. Then, for each $n$, we apply the triangle inequality to estimate
\begin{align}\label{e:Q-estimate}
    \Big| \QF(a', \xi')[\psi, \theta] - \QF(a_n^*, \xi_n^*)[\psi, \theta] \Big|
    \leq&
    \phantom{+\;}
    \abs{  \avg{a' |\nabla^{\perp} \psi|^2 
        -  a_n^* |\nabla^{\perp} \psi_n|^2} }
    \\ \nonumber
    &+ \abs{ \avg{ \abs{\nabla \theta}^2 - \abs{\nabla \theta_n}^2}}
    \\ \nonumber
    &+ \abs{ \avg{ \xi' \,\nabla^{\perp} \psi \cdot \nabla \theta
            - \xi^*_n \,\nabla^{\perp} \psi_n \cdot \nabla \theta_n  } } \Pe.
\end{align}
The first two terms on the right-hand side tend to zero with increasing $n$ because $a_n^* \to a'$, $\psi_n \to \psi$ strongly, and $\theta_n \to \theta$ strongly. The last term also tends to zero, but showing this requires more care because the convergence $\xi_n^* \weakto \xi'$ is only weak. (We cannot rely on strong convergence because we derived strong convergence using \cref{e:claim2}, which is what we are proving.) To handle this difficulty, we use the triangle inequality and estimate
\begin{align*}
    \abs{ 
        \avg{ \xi' \nabla^{\perp} \psi \cdot \nabla \theta
        - \xi^*_n \nabla^{\perp} \psi_n \cdot \nabla \theta_n
        }
    }
    \leq & 
    \phantom{+\;}
    \abs{
        \avg{
            (\xi' - \xi_n^*) \nabla^{\perp} \psi \cdot \nabla \theta
        }
    }
    \\
    &+
    \abs{
        \avg{
            \xi^*_n \nabla^{\perp} (\psi - \psi_n) \cdot \nabla \theta
        }
    }
    \\ 
    &+ 
    \abs{
        \avg{
            \xi^*_n \nabla^{\perp}\psi_n \cdot \nabla (\theta - \theta_n)
        }
    }
    .
\end{align*}
We then apply estimate \cref{e:clms} to the last two terms in this inequality and obtain
\begin{align*}
    \abs{ 
        \avg{ \xi' \nabla^{\perp} \psi \cdot \nabla \theta
        - \xi^*_n \nabla^{\perp} \psi_n \cdot \nabla \theta_n
        }
    }
    \lesssim &
    \phantom{+\;}
    \abs{
        \avg{
            (\xi' - \xi_n^*) \nabla^{\perp} \psi \cdot \nabla \theta
        }
    }
    \\
    &+ 
    \| \nabla \xi_n^* \|_{L^2}
    \| \nabla \psi - \nabla\psi_n \|_{L^2} 
    \| \nabla \theta \|_{L^2}  
    \\
    & + 
    \| \nabla \xi_n^* \|_{L^2}
    \| \nabla \psi_n \|_{L^2} 
    \| \nabla \theta - \theta_n \|_{L^2}
    .
\end{align*}
The weak convergence $\xi_n^* \weakto \xi'$ in $\smash{\Hoz}$, the strong convergence $\psi_n\to \psi$ in $\smash{H^1(\Omega)}$, and the strong convergence $\theta_n \to \theta$ in $\smash{\Hoz}$ ensure that all three terms on the right-hand side of this inequality tend to zero. Thus, the last term on the right-hand side of \cref{e:Q-estimate} vanishes with increasing $n$. Since the first two terms in that inequality also vanish, we deduce that
\begin{equation*}
    \QF(a', \xi')[\psi, \theta] = 
    \lim_{n \to \infty}  \; \QF(a_n^*, \xi_n^*)[\psi_n, \theta_n]
    \geq 0,
\end{equation*}
where the last inequality holds because $a_n^*$ and $\xi_n^*$ are feasible for \cref{e:ldp-discrete}. We have thus proved \cref{e:claim2}, completing the proof of \cref{thm:convergence}.
\end{proof}
\section{Numerical demonstrations}\label{s:numerics}

We demonstrate the computational approach of \cref{s:discretization} on two examples. In the first, we optimize bounds on the optimal cost $\ocpval(\Pe)$ of problem \cref{e:ocp} for two heat distributions in the square $\Omega=\smash{(-1,1)^2}$. In the second, we consider a uniformly heated annulus. 
In both cases, we exploit symmetries in the domain's geometry and in the heat distribution to replace the dual problem \cref{e:ldp} with equivalent problems, whose SDP discretizations have a lower computational complexity. Our convergence analysis extends to these symmetry-reduced formulations with straightforward modifications. We implemented all SDPs in Julia using \texttt{JuMP}~\cite{jump} and the interior-point solver \texttt{Mosek}~\cite{mosek}.

%%%%%%%%%%%%%%%%%%%%%%%%%%%%%%%%%%%%%%%%%%%%%%%%%%%%%%%%
\subsection{Optimal cooling in a square}\label{ss:example-square}
For the square domain $\Omega = (-1,1)^2$, we consider two heat source distributions, 
$f(\vec{x})=1$ and $f(\vec{x}) = 1 + \cos(\frac\pi2 x_1) \cos(\frac{3\pi}{2} x_2)$. 
The streamfunction space is $\cal{H} = \Hoz$ because the domain is simply connected.

\subsubsection{Symmetry reduction}
\label{ss:example-square-symm}
Since the domain and the heat distributions we consider are invariant under the transformations $(x_1, x_2)\mapsto(-x_1, x_2)$ and $(x_1, x_2)\mapsto(x_1, -x_2)$, we can to simplify problem \cref{e:ldp} before discretizing it into an SDP. Specifically, we decompose $H^1_0(\Omega)$ into the direct sum of the subspaces
\begin{gather*}
    \Hpp = \left\{u \in H^1_0(\Omega): u(-x_1, x_2)=u(x_1, -x_2)=u(x_1, x_2)\right\},\\
    \Hpm = \left\{u \in H^1_0(\Omega): u(-x_1, x_2)=u(x_1, x_2),\, u(x_1, -x_2)=-u(x_1, x_2)\right\},\\
    \Hmp = \left\{u \in H^1_0(\Omega): u(-x_1, x_2)=-u(x_1, x_2),\, u(x_1, -x_2)=u(x_1, x_2)\right\},\\
    \Hmm = \left\{u \in H^1_0(\Omega): u(-x_1, x_2)=u(x_1, -x_2)=-u(x_1, x_2)\right\},
\end{gather*}
and we rewrite problem \cref{e:ldp} as
\begin{align}\label{e:ldp-square-symm}
    \lb(\Pe) = \max_{\substack{a \in \R \\ \xi \in \Hpp}}\quad
    &\int_\Omega f\xi - \frac14 \abs{\nabla \xi}^2 \,\dx - a
    \\
    \nonumber
    \text{s.t.}\quad
    &\QF\left(a,\xi\right)[\psi,\theta] \geq 0 \quad \forall (\psi,\theta) \in \Hpp \times \Hmm,\\\nonumber
    &\QF\left(a,\xi\right)[\psi,\theta] \geq 0 \quad \forall (\psi,\theta) \in \Hpm \times \Hmp,\\\nonumber
    &\QF\left(a,\xi\right)[\psi,\theta] \geq 0 \quad \forall (\psi,\theta) \in \Hmm \times \Hpp,\\\nonumber
    &\QF\left(a,\xi\right)[\psi,\theta] \geq 0 \quad \forall (\psi,\theta) \in \Hmp \times \Hpm.
\end{align}
Indeed, the restriction of the optimization to $\xi \in \Hpp$ is warranted by a standard symmetrization argument and the convexity of \cref{e:ldp}. Given this restriction, it suffices to check that $\QF\left(a,\xi\right)\succeq 0$ on the listed combinations of subspaces because, for all other combinations, the symmetries of $\psi$ and $\theta$ imply that $\QF\left(a,\xi\right)[\psi,\theta]=0$.

The last two constraints in \cref{e:ldp-square-symm} are actually redundant. To see why, recall that the optimal $a$ is strictly positive. Then, given $\psi \in \Hmm$ and $\theta \in \Hpp$, we have that $\QF(\psi,\theta)=\QF(\psi',\theta')$ for $\psi'=\theta / \sqrt{a}\in \Hpp$ and $\theta'=-\sqrt{a}\psi \in \Hmm$, so the first and third constraints are equivalent. The same is true for the second and fourth constraints. In our implementation, therefore, we consider only the first two constraints. This is an example of the computational savings mentioned in \cref{rem:symmetrization-2}.

\subsubsection{Discretization scheme}
To discretize problem \cref{e:ldp-square-symm}, we use a spectral method based on expansions with respect to the eigenfunctions of the Dirichlet Laplacian on $\Omega$, which are simply sinusoidal functions. Specifically, given a positive integer $n\in \N$, we replace the spaces $\Hpp$, $\Hpm$, $\Hmp$ and $\Hmm$ with their $n^2$-dimensional subspaces
\begin{gather*}
    \Hpp_n = \bigg\{
        \sum_{\alpha,\beta\leq n} 
            \hat{u}_{\alpha\beta}
            \cos\tfrac{(2\alpha-1) \pi x}{2}
            \cos\tfrac{(2\beta-1) \pi y}{2}
            :\; \hat{u}_{\alpha\beta} \in \R
        \bigg\}
    \\
    \Hpm_n = \bigg\{
        \sum_{\alpha,\beta\leq n} 
            \hat{u}_{\alpha\beta}
            \cos\tfrac{(2\alpha-1) \pi x}{2}
            \sin\left(\beta \pi y\right)
            :\; \hat{u}_{\alpha\beta} \in \R
        \bigg\}
    \\
    \Hmp_n = \bigg\{
        \sum_{\alpha,\beta\leq n} 
            \hat{u}_{\alpha\beta}
            \sin\left(\alpha \pi x\right)
            \cos\tfrac{(2\beta-1) \pi y}{2}
            :\; \hat{u}_{\alpha\beta} \in \R
        \bigg\}
    \\
    \Hmm_n = \bigg\{
        \sum_{\alpha,\beta\leq n} 
            \hat{u}_{\alpha\beta}
            \sin\left(\alpha \pi x\right)
            \sin\left(\beta \pi y\right)
            :\; \hat{u}_{\alpha\beta} \in \R
        \bigg\}
\end{gather*}
With these choices, problem \cref{e:ldp-square-symm} (without the last two redundant constraints) is approximated by an SDP with $n^2+1$ optimization variables and two LMIs of size $2n^2 \times 2n^2$.

\subsubsection{Results}
For each of the two heat distributions reported at the start of this section, we solved SDP discretizations of the symmetry-reduced problem \cref{e:ldp-square-symm} at $\Pe=1$, $10$ and $20$ with discretization parameter $n$ ranging from $1$ to $30$. The corresponding lower bounds $\lb_n(\Pe)$ are reported in \cref{tab:square-results} and converge quickly as $n$ is increased. One would expect such a fast convergence from a spectral discretization, even though the analysis in \cref{s:discretization} does not provide any convergence rates. It is also clear from the table that higher values of $\Pe$ necessitate a higher resolution for a fixed level of accuracy.

\begin{table}[t]
    \centering
    \caption{
        \label{tab:square-results}
        Lower bounds $\lb_n(\Pe)$ for optimal cooling in a square, computed at $\Pe=1$,$10$ and $20$ for two heat distributions $f$ and increasing values of the discretization parameter $n$.
    }
    \begin{tabular}{c c ccc c ccc}
        \toprule
        && \multicolumn{3}{c}{$f(\vec{x}) = 1$}
        && \multicolumn{3}{c}{$f(\vec{x}) = 1 + \cos(\frac\pi2 x_1) \cos(\frac{3\pi}{2} x_2)$}
        \\
        \cmidrule{3-5} \cmidrule{7-9}
        $n$
        && $\lb_n(1)$  & $\lb_n(10)$  & $\lb_n(20)$ 
        && $\lb_n(1)$  & $\lb_n(10)$  & $\lb_n(20)$ 
        \\[0.75ex]
          1 && 0.530145 & 0.365785 & 0.188599 && 0.530145 & 0.365785 & 0.188599\\ 
          2 && 0.554924 & 0.424603 & 0.270803 && 0.562662 & 0.425173 & 0.261577\\ 
          3 && 0.558572 & 0.441014 & 0.300804 && 0.568036 & 0.433713 & 0.290004\\ 
          4 && 0.559561 & 0.443606 & 0.314503 && 0.569504 & 0.436696 & 0.304565\\ 
          5 && 0.559932 & 0.444797 & 0.319255 && 0.570054 & 0.437968 & 0.309230\\[0.75ex] % easier to read
         10 && 0.560286 & 0.445523 & 0.322775 && 0.570578 & 0.438676 & 0.312156\\ 
         15 && 0.560323 & 0.445563 & 0.322890 && 0.570633 & 0.438733 & 0.312266\\ 
         20 && 0.560333 & 0.445572 & 0.322903 && 0.570646 & 0.438746 & 0.312282\\ 
         25 && 0.560336 & 0.445576 & 0.322907 && 0.570651 & 0.438751 & 0.312287\\ 
         30 && 0.560337 & 0.445577 & 0.322908 && 0.570653 & 0.438753 & 0.312289\\ 
        \bottomrule
    \end{tabular}
\end{table}

\begin{figure}
    \centering
    \includegraphics[scale=0.92]{img/square_f2_Pe_10_n30.pdf}
    \caption{
    \label{fig:square-sin-critical-fields-Pe10}
    Results for optimal cooling in a square with a sinusoidal heat distribution, obtained at $Pe=10$ with discretization parameter $n=30$.
    \emph{Top-left:} The heat distribution $f$.
    \emph{Bottom-left:} The optimal field $\xi^*$ for \cref{e:ldp-square-symm}.
    \emph{Top-right and bottom-right:} Critical pairs $(\psi_i, \theta_i)$, for which the first two constraints $\QF(\psi,\theta)\geq 0$ in \cref{e:ldp-square-symm} are active.
    }
\end{figure}
\begin{figure}
    \centering
    \includegraphics[scale=0.92]{img/square_f2_Pe_20_n30.pdf}
    \caption{
    \label{fig:square-sin-critical-fields-Pe20}
    Results for optimal cooling in a square with a sinusoidal heat distribution, obtained at $Pe=20$ with discretization parameter $n=30$.
    \emph{Top-left:} The heat distribution $f$.
    \emph{Bottom-left:} The optimal field $\xi^*$ for \cref{e:ldp-square-symm}.
    \emph{Top-right and bottom-right:} Critical streamfunction-temperature pairs $(\psi_i, \theta_i)$, for which the first two constraints $\QF(\psi,\theta)\geq 0$ in \cref{e:ldp-square-symm} are active.
    }
\end{figure}

The reason why a finer discretization is needed as $\Pe$ is increased becomes clear if one inspects the optimal solution of problem \cref{e:ldp-square-symm}. For brevity, we focus only on the sinusoidal heat distribution at $\Pe=10$ and $\Pe=20$; results for the uniform heating case are similar. \Cref{fig:square-sin-critical-fields-Pe10,fig:square-sin-critical-fields-Pe20} show that the optimal solution $\xi^*$ has a relatively simple structure at both $\Pe$ values. In contrast, the pairs $(\psi_i,\theta_i)$ for which the first two constraints in \cref{e:ldp-square-symm} are active become increasingly complex. (Additional pairs for which the other two redundant constraints are active can be recovered using the symmetry transformation described at the end of \cref{ss:example-square-symm} and are not shown for brevity). Accurately resolving these fields is necessary to obtain a well-converged solution. Note also that the number of such `critical fields' increases with $\Pe$, in analogy to what is observed when implementing the background method for estimating time-average properties of turbulent flows between parallel plates (see, e.g., \cite{Doering1996,Plasting2003,fpw2018,Wen2015}). 

%%%%%%%%%%%%%%%%%%%%%%%%%%%%%%%%%%%%%%%%%%%%%%%%%%%%%%%%%%%%%%%%
\subsection{Optimal cooling in a uniformly heated annulus}\label{ss:annulus}
Next, we consider optimal cooling in the annulus $\Omega = \{\vec{x}\in\R^2:\;\rho < \abs{\vec{x}} < 1\}$ with $\rho=\frac14$ and a uniform heat distribution $f(\vec{x})=1$. It is convenient to work in the polar coordinates $(r,\phi)$.

\subsubsection{Symmetry reduction}\label{ss:annulus-symmetrization}
Since the chosen heat distributions depend only on the radial direction, we may restrict the optimization to functions $\xi=\xi(r)$ varying only in the radial direction. Then, a Fourier series expansion in polar coordinates detailed in \cref{app:fourier-annulus} allows us to rewrite problem \cref{e:ldp} as
\begin{equation}\label{e:ldp-annulus-symm}
    \lb(\Pe) = 
    \max_{\substack{a \in \R \\ \xi \in H^1_0(\rho,1) \\ \QF_k(a,\xi)\succeq 0 \,\forall k\in\Z}}\quad
    2\pi\int_\rho^1  \left( f\xi - \frac{1}{4} \abs{\xi'}^2 \right) r\dr - a,
\end{equation}
where $k$ is the wavenumber of each Fourier mode in the expansion and $\QF_k(a,\xi)$ is a quadratic form on $H^1_0(\rho,1) \times H^1_0(\rho,1)$ defined as
\begin{equation}\label{e:qf-fourier-annulus}
    \QF_k(a,\xi)[\psi,\theta] = \int_\rho^1 
        \left[ a \left(r|\psi'|^2 + \frac{k^2}{r}|\psi|^2\right) +
         \left( r|\theta'|^2 + \frac{k^2}{r} |\theta|^2 \right)
         + k\Pe\,\xi'\,\psi \theta  \right] \dr.
\end{equation}

Problem~\cref{e:ldp-annulus-symm} can be simplified further by setting $a=\alpha^2$, $\psi=(u+v)/\sqrt{\alpha}$ and $\theta = (u-v)\sqrt{\alpha}$, which is a version of the change of variables discussed in \cref{rem:symmetrization-1}. Indeed, a direct calculation shows that $\QF_k(a,\xi)[\psi,\theta] = 2\mathcal{S}_k^+(a,\xi)[u] + 2\mathcal{S}_k^-(a,\xi)[v]$ where
\begin{align*}
    S_k^\pm(a,\xi)[u] = \int_\rho^1 
        \alpha \left(r|u'|^2+ \frac{k^2}{r}|u|^2 \right)
        \pm \frac{\Pe}2 \,k\, \xi' u^2  \, \dr.
\end{align*}
Since the functions $u$ and $v$ can be chosen independently, $\QF_k(a,\xi) \succeq 0$ on $H^1_0(\rho,1) \times H^1_0(\rho,1)$ for every $k\in\Z$ if and only if $\mathcal{S}_k^+(a,\xi) \succeq 0$ on $H^1_0(\rho,1)$ for every $k\in\Z$. Consequently, the maximization problem in \cref{e:ldp-annulus-symm} is equivalent to
\begin{align}\label{e:ldp-annulus-symm-v2}
    \max_{\substack{\alpha \in \R \\ \xi \in H^1_0(\rho,1) \\ S_k^+(a,\xi) \succeq 0 \,\forall k \in \Z}}\quad
    &2\pi\int_\rho^1  \left( f\xi - \frac{1}{4} \abs{\xi'}^2 \right) r\dr - \alpha^2.
\end{align}


\subsubsection{Discretization scheme}
To discretize problem \cref{e:ldp-annulus-symm-v2}, we first fix $m\in\N$ and impose the constraint $S_k^+(a,\xi) \succeq 0$ only for nonzero $k \in \{-m,\ldots,m\}$. (The constraint is trivially satisfied if $k=0$.) We then discretize the radial direction using a pseudo-spectral Chebyshev collocation method with $n+2$ grid points, that is, we replace $H^1_0(\rho,1)$ with the space of polynomials of degree $n+1$ defined by their values at $n+2$ Chebyshev points of the second kind. We compute the integrals in \cref{e:ldp-annulus-symm-v2} and in the definition of $S_k^+(a,\xi)$ using a Clenshaw--Curtis quadrature at $3(n+2)$ Chebyshev points of the second kind. This quadrature is exact for the polynomials $f\xi$, $|\xi'|^2$, $r|u'|^2$, and $\xi'u^2$, and is also accurate for the rational function $u^2/r$. 

Upon accounting for the homogeneous Dirichlet boundary conditions satisfied by functions in $H^1_0(\rho,1)$, this discretization strategy results in an SDP with $n+1$ optimization variables and $2m$ LMIs of size $n\times n$ each. We consider $m$ and $n$ sufficiently large if halving them changes the numerical approximation of $\lb(\Pe)$ by less than $0.1\%$. With this criterion, the largest \Peclet\ number we could consider on a workstation with $500$GB RAM before exceeding the available memory was $\Pe = 354.8$, which required $(m,n)=(256,64)$.

To access higher \Peclet\ numbers, we used an alternative approach where we impose the constraint $S_k^+(a,\xi) \succeq 0$ for all $k\in\R$. After discretization in the radial direction, this stronger constraint becomes an $n\times n$ quadratic polynomial matrix inequality with indeterminate $k$, and can be enforced computationally using techniques for sum-of-squares programming (see, e.g., \cite[\S3.3.2]{Parrilo2013}). The resulting optimization problem is an SDP with one positive semidefinite cone of dimension $2n\times 2n$, which we implement using the Julia package \texttt{SumOfSquares.jl} \cite{SumOfSquares1,SumOfSquares2}. Again, we consider $n$ sufficiently large if halving it changes the SDP's optimal value by less than $0.1\%$. This strategy allowed us to access \Peclet\ numbers up to $\Pe=1000$, which required $n=256$ and (at this resolution) approximately 15 hours per SDP.

\subsubsection{Results}
We used the two discretization schemes described above to approximate $\lb(\Pe)$ for \Peclet\ numbers ranging from $\Pe=0.1$ to $\Pe=1000$. The results are plotted in the top-left panel of \cref{f:bounds_vs_Pe_annulus}. Other panels in that figure show the optimal constant $a$ as a function of $\Pe$ and the optimal function $\xi(r)$ for $\Pe=10$, $100$ and $1000$. Note that the results obtained with the discrete wavenumbers $k\in\Z$ are indistinguishable from those obtained with $k\in\R$, so we only discuss the latter.

\begin{figure}[t]
    \centering
    \includegraphics[width=\linewidth]{img/LB_vs_Pe_annulus.pdf}
    \caption{Optimal solutions of SDP discretizations of problem \cref{e:ldp-annulus-symm-v2}, which gives lower bounds on the optimal value $\ocpval(\Pe)$ of problem \cref{e:ocp} for a uniformly heated annulus. Results are computed by imposing the constraint $S_k^+(a,\xi) \succeq 0$ for $k\in\{-m,\ldots,m\}$ (red lines) and for all $k\in\R$ (blue dots). \emph{Top-left:} Lower bounds as a function of $\Pe$. \emph{Bottom-left:} Optimal $a$ as a function of $\Pe$. \emph{Right column:} Optimal $\xi(r)$ for $\Pe=10$, $100$ and $1000$. Note the different ranges of the vertical axis.
    \label{f:bounds_vs_Pe_annulus}
    }
\end{figure}

Our numerical approximations to $\lb(\Pe)$ are consistent with the functional form of the analytical lower bound from \cref{th:analytical-lb}, namely, $\lb(\Pe) \sim 1/[1 + g(\Pe)]$ for some increasing and unbounded function $g(\Pe)$. However, we cannot determine the asymptotic decay rate of $\lb(\Pe)$ with any confidence because the range of \Peclet\ numbers spanned by our computations is too small.  
\begin{figure}
    \centering
    \includegraphics[width=0.8\linewidth]{img/critical_k_annulus.pdf}
    \caption{
    \emph{Left:} Active wavenumbers $k$ for problem \cref{e:ldp-annulus-symm-v2} as a function of $\Pe$. \emph{Right:} The same plot rescaled by $\Pe$, illustrating the superlinear growth of the largest active wavenumber. Both plots are obtained from discretizations of \cref{e:ldp-annulus-symm-v2} that consider $k\in\R$; similar plots for $k\in\Z$ are omitted for brevity.}
    \label{f:annulus-critical-k}
\end{figure}

The main barrier to accessing a broader range of \Peclet\ numbers is the high resolution needed to make our SDP discretizations accurate, which is due to two factors. First, although the constraint $S_k^+(a,\xi) \succeq 0$ is active only for a small set of wavenumbers $k$ (cf. \cref{f:annulus-critical-k}), the active set cannot be determined \emph{a priori}, so SDP discretizations must include a large number of LMIs. Second, even if one reduces this number using an `active set' method or an iterative wavenumber selection strategy similar to that in \cite[\S5.2]{fpw2018}, the largest active wavenumbers appear to increase super-linearly with $\Pe$. Approximating the corresponding critical functions $u$ requires a high resolution because they correspond to flows with a high-frequency roll structure localized near the domain boundaries, as shown in \cref{f:annulus-critical-flows} for the moderate $\Pe=100$.

On the other hand, it is clear from our results that the optimal $\xi$ obtained computationally differs significantly from that used in the proof of \cref{th:analytical-lb}, which for the present annular geometry is a $\Pe$-dependent multiple of $\smash{(\frac{1-\rho^2}{\ln \rho}) \ln r + r^2 - 1}$. In the next section, we will demonstrate that analytical choices of $\xi$ inspired by the profiles plotted in \cref{f:bounds_vs_Pe_annulus} can improve the efficiency bounds in \cref{th:analytical-lb} by a logarithmic factor.

\begin{figure}[t]
    \centering
    \includegraphics[scale=1]{img/critical_modes_Pe100a.pdf}\\[10pt]
    \includegraphics[scale=1]{img/critical_modes_Pe100b.pdf}
    \caption{Critical streamfunction-temperature pairs $(\psi,\theta)$ and wavenumbers $k$ for which the constraint of problem \cref{e:ldp-annulus-symm} are active at $\Pe=100$. These pairs are recovered from functions $u(r)$, also shown, for which the constraint of problem \cref{e:ldp-annulus-symm-v2} is active.
    \label{f:annulus-critical-flows}
    }
\end{figure}
\section{Improved analytical bounds for the annulus and the disk}
\label{s:improved-bounds}

Inspired by the numerical results presented in \cref{ss:annulus}, we now take a closer look at optimal cooling problems in circular and annular domains. By a suitable choice of coordinates, it suffices to consider the unit annulus $\Omega=\{\vec{x}\in\R^2: \rho < \abs{\vec{x}} < 1\}$ with inner radius $\rho$ and the unit disk $\Omega=\{\vec{x}\in\R^2: \abs{\vec{x}} < 1\}$. We will show that the optimal cooling efficiency bound $\efficiency(\Pe) \lesssim \Pe^{2}$ from \cref{th:analytical-lb} can be reduced by a logarithmic factor if the heat distribution $f$ has a strictly positive average on any circle centered at the origin. This class includes, for example, the uniform heat distribution $f=1$ studied in \cref{ss:annulus} and all positive and radially symmetric $f$.


Specifically, let $f$ be given in the polar coordinates $(r,\phi)$ for convenience, and let
\begin{equation*}
    \overline{f}(r) := \frac{1}{2\pi}
    \int_0^{2\pi} f(r,\phi) \, {\rm d}
    \phi
\end{equation*}
denote its azimuthal average.
We prove the following results.

\begin{theorem}[Bounds for the annulus]
\label{th:refined-annulus}
    Let $\Omega=\{\vec{x}\in\R^2: \rho < \abs{\vec{x}} < 1\}$ be a unit annulus with inner radius $\rho \in(0,1)$. 
    Assume $\overline{f}(r) \geq f_{\rm{min}} > 0$ for all $r \in (\rho,1)$. There exists a constant $C(\rho,f) \geq 16 \pi^2 \rho^6 (1-\rho)^3 f_{\rm min}^2 / (1-\rho + 8\pi\rho^4 \e)$ such that
    \begin{equation*}
        \ocpval(\Pe) \geq 
        C(\rho,f_{\rm min}) \frac{\ln^2 \Pe}{\Pe^2} \qquad \forall \, \Pe>0.
    \end{equation*}
    Consequently, there exists a constant $C'(\rho,f)$ such that $\efficiency(\Pe) \leq C'(\rho,f) \frac{\Pe^2}{\ln^2 \Pe}$.
\end{theorem}

\begin{theorem}[Bounds for the disk]
\label{th:refined-disk}
    Let $\Omega$ be the unit disk. Assume $\overline{f}(r) \geq f_{\rm{min}} > 0$ for all $r \in [0,1)$. There exists a constant $C(f_{\rm min}) \geq \pi^2 f_{\rm min}^2 / (64 + 32 \pi \e)$ such that
    \begin{equation*}
        \ocpval(\Pe) \geq 
        C(f_{\rm min}) \frac{\ln^2 \Pe}{\Pe^2} \qquad \forall \, \Pe>0.
    \end{equation*}
    Consequently, there exists a constant $C'(f)$ such that $\efficiency(\Pe) \leq C'(f) \frac{\Pe^2}{\ln^2 \Pe}$.
\end{theorem}

\Cref{th:refined-annulus} is proved in \cref{ss:proof-annulus}. The proof of \cref{th:refined-disk} is very similar and relies on the same key estimates, so we only outline it in \cref{ss:proof-disk}.


\subsection{Proof of \texorpdfstring{\cref{th:refined-annulus}}{Theorem~\ref{th:refined-annulus}}}\label{ss:proof-annulus}

We only need to prove the lower bound on $\ocpval(\Pe)$, since the corresponding upper bound on $\efficiency(\Pe)$ follows immediately from \cref{e:efficiency}.
Restricting the maximization in the dual problem \cref{e:ldp} to functions $\xi$ that depend only on the radial position $r=|\vec{x}|$ gives the lower bound
\begin{equation}\label{e:annulus-general-f}
    \ocpval(\Pe) \geq 
    \lb(\Pe) \geq \max_{\substack{a \in \R\\ \xi \in H^1_0(\rho,1) \\ \QF(a,\xi)\succeq 0}} \;2\pi\int_{\rho}^1 \left( \xi \overline{f} - \frac14 |\xi'|^2 \right) r \dr - a. 
\end{equation}
We can then apply the same symmetry reduction arguments outlined in \cref{ss:annulus-symmetrization} to conclude that the maximization problem on the right-hand side of this inequality can be rewritten as \cref{e:ldp-annulus-symm-v2} with $\overline{f}$ in place of $f$. We will construct a feasible $\xi$ and $\alpha = \sqrt{a}$ for that problem, which achieve an objective value no smaller than a multiple of $\smash{\frac{\ln^2 \Pe}{\Pe^2}}$.

\subsubsection{The candidate feasible point}
Let $\rho_0= \frac12(1+\rho)$ denote the average radius of the annulus. We claim that suitable choices of $\alpha$ and $\xi$ are
\begin{equation}\label{e:annulus-proof-xi-alpha}
    \alpha
    =
    \frac{\delta}{2\rho^2} \, \Pe
    \qquad\text{and}\qquad
    \xi(r) = \begin{cases}
        \delta \ln\left( \frac{r-\rho + \varepsilon}{\varepsilon}\right) & \text{if } r \in (\rho, \rho_0], \\
        \delta \ln\left( \frac{1 + \varepsilon - r}{\varepsilon}\right) & \text{if }r \in (\rho_0, 1),
    \end{cases}
\end{equation}
where $\delta$ and $\varepsilon$ are positive $\Pe$-dependent constants given by
\begin{equation*}
    \varepsilon = \frac{1-\rho}{2 \e}\,\frac{1}{\Pe^2}
    \qquad\text{and}\qquad
    \delta
    =
    \frac{8 \pi f_{\rm min} \rho^5 (1-\rho)^2}{1-\rho + 8\pi \e \rho^4} \,
    \frac{\ln\Pe}{\Pe^2}.
\end{equation*}
\Cref{f:analytical-xi} shows the function $\xi$ for $\Pe=20$, $\rho=\frac14$ and $f_{\rm min} = 1$. Note the qualitative similarity with the optimal $\xi$ obtained numerically at large $\Pe$ for $f=1$, shown in \cref{f:bounds_vs_Pe_annulus}.

\begin{figure}
    \centering
    \includegraphics[scale=1]{img/analytical_xi.pdf}
    \caption{Analytical $\xi$ used in the proof of \cref{th:refined-annulus} (left) and in the proof of \cref{th:refined-disk} (right). The plots are for $\Pe=20$, $\rho=\frac14$ and $f_{\rm min} = 1$.}
    \label{f:analytical-xi}
\end{figure}

\subsubsection{Feasibility} 
We now verify that the function $\xi$ and the constant $\alpha$ in \cref{e:annulus-proof-xi-alpha} satisfy the constraints of problem \cref{e:ldp-annulus-symm-v2}, so they are admissible for the maximization problem in~\cref{e:annulus-general-f}. First, we establish a version of Hardy's inequality.

\begin{lemma}\label{lemma:hardy}
    Every $u\in H_{0}^{1}(\rho,1)$ satisfies
    \begin{align*}
        \int_{\rho}^{\rho_0}\frac{1}{(r-\rho)^{2}}\,\frac{u^{2}}{r}\,\dr	
        &\leq\frac{4}{\rho^{2}}\int_{\rho}^{\rho_0}r\abs{u'}^{2}\,\dr,
        \\
        \int_{\rho_0}^{1}\frac{1}{(1-r)^{2}}\,\frac{u^{2}}{r}\,\dr	
        &\leq\frac{4}{\rho^{2}}\int_{\rho_0}^{1}r\abs{u'}^{2}\,\dr.
    \end{align*}
\end{lemma}

\begin{proof}
    Since $u(\rho)=0$, we can use Hardy's inequality \cite[Exercise~8.8]{Brezis2011} to estimate 
    \begin{eqnarray*}
    \int_{\rho}^{\rho_0}\frac{1}{(r-\rho)^{2}}\frac{u^{2}}{r}\,\dr
    &\leq &
    \frac{1}{\rho} \int_{\rho}^{\rho_0}\frac{u^{2}}{(r-\rho)^{2}}\,\dr
    \\
    &\overset{\substack{\text{Hardy's}\\\text{ineq.}\\[0.5ex]}}{\leq}& 
    \frac{4}{\rho} \int_{\rho}^{\rho_0}\abs{u'}^{2}\,\dr
    \\
    &\leq &
    \frac{4}{\rho^{2}} \int_{\rho}^{\rho_0}r\abs{u'}^{2}\,\dr.
    \end{eqnarray*}
    This proves the first inequality. The second one follows similarly.
\end{proof}

Next, we use these Hardy-type estimates to establish the claimed feasibility.

\begin{proposition}\label{prop:annulus-feasibility}
    Let $\xi$ and $\alpha$ be given as in \cref{e:annulus-proof-xi-alpha}. For all $u\in H^1_0(\rho,1)$ and all $k \in \Z$,
    \begin{equation*}
        \int_\rho^1 
        \alpha \left(r\abs{u'}^2+ \frac{k^2}{r}u^2 \right)
        + \frac{\Pe}2 \,k\, \xi' u^2 \, \dr \geq 0.
    \end{equation*}
\end{proposition}

\begin{proof}
    Since $\alpha= \Pe\,\delta / (2\rho^2)$ by definition, and since
    \begin{align*}
        \left|\int_{\rho}^{1}\frac{k}{2}\Pe\, \xi'u^{2}\,\dr\right| 
        \leq
        \frac{\Pe\, \delta}{2} \int_{\rho}^{\rho_0}\frac{ku^{2}}{r-\rho+\varepsilon}\,\dr
        +\frac{\Pe \,\delta}{2}
        \int_{\rho_0}^{1}\frac{ku^{2}}{1+\varepsilon-r}\,\dr,
    \end{align*}
    the claim follows if we can show that
    \begin{subequations}
        \begin{align}
        \label{e:annulusc-cnstr-1}
        \int_{\rho}^{\rho_0}\frac{ku^{2}}{r-\rho+\varepsilon}\,\dr 
        &\leq 
        \frac{1}{\rho^{2}} \int_{\rho}^{\rho_0}\left(r\abs{u'}^{2}+k^{2}\right)\frac{u^{2}}{r}\,\dr,
        \\
        \label{e:annulusc-cnstr-2}
        \int_{\rho_0}^{1}
        \frac{ku^{2}}{1+\varepsilon-r}\,\dr 
        &\leq 
        \frac{1}{\rho^{2}} \int_{\rho_0}^{1}
        \left(r\abs{u'}^{2}+k^{2}\right)\frac{u^{2}}{r}\,\dr.
    \end{align}
    \end{subequations}
    To derive \cref{e:annulusc-cnstr-2}, we use the inequalities $\smash{\frac{1}{1+\varepsilon - r}}< \smash{\frac{1}{1-r}}$ and $1 < \smash{\frac1r}$ on $[\rho_0,1)$, the inequality $ab \leq \smash{\frac14}a^2+b^2$, the second inequality in \cref{lemma:hardy}, and the inequality $1<\smash{\frac{1}{\rho^2}}$ to estimate
    \begin{align*}
        \int_{\rho_0}^1 
        \frac{ku^{2}}{1+\varepsilon-r}\,\dr 
        &\leq 
        \int_{\rho_0}^1\frac{k}{1-r}\,u^{2}\,\dr
        \\
        &\leq
        \int_{\rho_0}^1 
        \left(\frac{1}{4(1-r)^{2}}+k^{2}\right)u^{2}\,\dr
        \\
        &\leq
       \int_{\rho_0}^1 
       \left(\frac{1}{4(1-r)^{2}}+k^{2}\right)\frac{u^{2}}{r}\,\dr
        \\
        &\leq
        \int_{\rho_0}^1
        \left(\frac{r}{\rho^{2}}\abs{u'}^{2}+\frac{k^{2}}{r}u^{2}\right)\,\dr
        \\
        &\leq
        \frac{1}{\rho^{2}}
        \int_{\rho_0}^1
        \left(r\abs{u'}^{2}+\frac{k^{2}}{r}u^{2}\right)\,\dr.
    \end{align*}
    Inequality \cref{e:annulusc-cnstr-1} follows analogously using the first inequality in \cref{lemma:hardy}.
\end{proof}


\subsubsection{Estimating the objective value}
\label{ss:annulus-proof-objective}

Having verified that the function $\xi$ and the constant $\alpha$ in \cref{e:annulus-proof-xi-alpha} are admissible for the maximization problem in \cref{e:annulus-general-f}, we now show that
\begin{equation}\label{e:annulus-proof-cost-bound}
    2\pi\int_{\rho}^1 \left( \xi \overline{f} - \frac14 |\xi'|^2 \right) r \dr - \alpha^2 \geq 
    C(\rho, f_{\rm min})
    \frac{\ln^{2}Pe}{Pe^{2}}
\end{equation}
for a suitable positive constant $C(\rho, f_{\rm min})$. This will conclude the proof of \cref{th:refined-annulus}.

Since $\xi$ is nonnegative and symmetric with respect to $r=\rho_0$ (see the right panel in \cref{f:analytical-xi}), we can estimate
\begin{align*}
    \int_{\rho}^1 \xi \overline{f}\, r\dr
    & \geq 
    2  f_{\rm min} \, \rho
    \int_{\rho_0}^1 \xi \dr
    \\
    &=
    2f_{\rm min} \, \delta \rho \left[ (1+\varepsilon-\rho_0) \ln\left( \frac{1+\varepsilon-\rho_0}{\varepsilon} \right) - (1-\rho_0)\right]
    \\
    &\geq 
    2f_{\rm min} \, \delta \rho (1-\rho_0) \left[ \ln\left( \frac{1-\rho_0}{\varepsilon} \right)  - 1\right]
    \\
    &=
    f_{\rm min} \, \delta \rho (1 - \rho) 
     \ln\left(\frac{1-\rho}{2 \e \varepsilon}\right).
\end{align*}
We also have
\begin{align*}
    \int_{\rho}^{1}|\xi'|^{2}r\,\dr 
    &\leq
    \int_{\rho}^{\rho_0}|\xi'|^{2}\,\dr+\int_{\rho_0}^{1}|\xi'|^{2}\,\dr
    \\
    &=
    \int_{\rho}^{\rho_0}\frac{\delta^{2}}{(r-\rho+\varepsilon)^{2}}\,\dr
    +
    \int_{\rho_0}^{1}\frac{\delta^{2}}{(1+\varepsilon-r)^{2}}\,\dr
    \\
    &=
    \delta^{2}\left(\frac{1}{\varepsilon}-\frac{1}{\rho_0-\rho+\varepsilon}+\frac{1}{\varepsilon}-\frac{1}{1+\varepsilon-\rho_0}\right)
    \\
    &
    \leq
    \frac{2\delta^{2}}{\varepsilon}.
\end{align*}
Using these inequalities, as well as the definitions of $\alpha$, $\delta$, and $\varepsilon$, we arrive (after some lengthy but straightforward algebra omitted for brevity) at
\begin{align*}
    2\pi\int_{\rho}^1 \left( \xi \overline{f} - \frac14 |\xi'|^2 \right) r \dr - \alpha^2 
    &\geq 
    2\pi f_{\rm min} \, \delta \rho (1 - \rho) 
    \ln\left(\frac{1-\rho}{2 \e \varepsilon}\right)
    -
    \frac{\pi \delta^{2}}{\varepsilon}
    - \alpha^2
    \\
    &= \frac{16 \pi^2 \rho^6 (1-\rho)^3  f_{\rm min}^2}{(1-\rho) + 8\pi\rho^4 \e}
    \,
    \frac{\ln^2\Pe}{\Pe^2}.
\end{align*}
This is precisely \cref{e:annulus-proof-cost-bound} with an explicit lower estimate for the constant $C(\rho,f_{\rm min})$.


\subsection{Proof of \texorpdfstring{\cref{th:refined-disk}}{Theorem~\ref{th:refined-disk}}}\label{ss:proof-disk}

We only need to prove the lower bound on $\ocpval(\Pe)$ since \cref{e:efficiency} implies the corresponding upper bound on $\efficiency(\Pe)$. Fix an arbitrary $\rho_0 \in (0,1)$. Set
\begin{equation*}
    \varepsilon = \frac{1-\rho_0}{2 \e} \, \frac{1}{\Pe^2}
    \qquad\text{and}\qquad
    \delta
    =
    \frac{4 \pi f_{\rm min} \rho_0^5 (1-\rho_0)^2}{1-\rho_0 + 4\pi \e \rho_0^4} \,
    \frac{\ln\Pe}{\Pe^2}.
\end{equation*}
We claim that
\begin{equation*}
    \alpha = \frac{\delta}{2\rho_0^2} \, \Pe
    \qquad\text{and}\qquad
    \xi(\vec{x}) = \begin{cases}
        \delta \ln\left( \frac{1 + \varepsilon - \rho_0}{\varepsilon}\right) & \text{if } \abs{\vec{x}} \in (0,\rho_0] \\
        \delta \ln\left( \frac{1 + \varepsilon - \abs{\vec{x}}}{\varepsilon}\right) & \text{if }\abs{\vec{x}} \in (\rho_0, 1)
    \end{cases}
\end{equation*}
are admissible for the `balanced' version of problem \cref{e:ldp} described in \cref{rem:symmetrization-1} and attain an objective value no smaller than a multiple of $\smash{\frac{\ln^2 \Pe}{\Pe^2}}$. (For illustration, the function $\xi$ is plotted in the right panel of \cref{f:analytical-xi} for $\Pe=20$, $f_{\rm min}=1$, and $\rho_0 = 0$.)

To verify this claim, let $B_{\rho_0}$ denote the disk of radius $\rho_0$ centered at the origin. Suppose for the time being that $\xi$ and $\alpha$ are indeed admissible. Since $\xi$ depends only on the radial coordinate $r=\abs{\vec{x}}$ and is a positive constant in $B_{\rho_0}$,  the assumption on $f$ in \cref{th:refined-disk} implies that $\int_{B_{\rho_0}} f \xi \dx > 0$. Then, the objective value attained by $\xi$ and $\alpha$ satisfies
\begin{align*}
    \avg{f \xi - \frac14\abs{\nabla \xi}^2} - \alpha^2
    &\geq 
    \int_{\Omega \setminus B_{\rho_0}} {f \xi - \frac14\abs{\nabla \xi}^2} \,\dx - \alpha^2
    \\
    &=2\pi \int_{\rho_0}^1 {\overline{f} \xi - \frac14\abs{\xi'}^2} \,r\dr - \alpha^2.
\end{align*}
The last expression can be bounded below with the same estimates used for the annular domain in \cref{ss:annulus-proof-objective}. We omit the details for brevity, but one eventually finds that
\begin{equation*}
    \avg{f \xi - \frac14\abs{\nabla \xi}^2} - \alpha^2
    \geq \frac{4\pi^2 \rho_0^6 (1-\rho_0)^3  f_{\min}^2}{1-\rho_0 + 4\pi \rho_0^4 \e} \, \frac{\ln^2 \Pe}{\Pe^2}.
\end{equation*}
The expression on the right-hand side is maximized when $\rho_0 \approx 0.4973$, so we set $\rho_0=\frac12$ to obtain the lower bound stated in \cref{th:refined-disk} with $C(f_{\rm min}) \geq \pi^2 f_{\rm min}^2 / (64 + 32 \pi \e)$.

It remains to verify that our choices of $\xi$ and $\alpha$ are admissible, meaning that the symmetrized quadratic form $\QF'(\alpha,\xi)[\psi,\theta]$ in \cref{e:qf-symmetrized} is nonnegative for all $\theta\in\Hoz$ and all $\psi \in \cal{H}$. Here, $\cal{H}=\Hoz$ because the disk $\Omega$ is simply connected. After integrating the $\xi$-dependent term in $\QF'(\alpha,\xi)[\psi,\theta]$ by parts and observing that $\xi$ is constant on $B_{\rho_0}$, we find that
\begin{align*}
    \QF'(\alpha,\xi)[\psi,\theta]
    &= 
    \avg{ \alpha |\nabla \psi|^2 + \alpha \abs{\nabla\theta}^2 + \Pe\,\theta \, \nabla^\perp \psi \cdot \nabla \xi}
    \\
    &\geq 
    \int_{\Omega \setminus B_{\rho_0}} { \alpha|\nabla \psi|^2 + \alpha \abs{\nabla\theta}^2 + \Pe\,\theta \, \nabla^\perp \psi \cdot \nabla \xi} \,\dx.
\end{align*}
Since the last integral is over the annulus $\Omega \setminus B_{\rho_0}$, we can analyze it with the same steps used for the annulus in \cref{ss:proof-annulus}, the only difference being a lack of boundary conditions on $\psi$ and $\theta$ when $\abs{\vec{x}}=\rho_0$. Specifically, we can perform the same Fourier series expansion described in \cref{app:fourier-annulus}, apply the symmetrization steps outlined in \cref{ss:annulus-symmetrization}, and conclude that $\QF'(\alpha,\xi)$ is positive semidefinite if
\begin{equation*}
    \int_{\rho_0}^1 
        \alpha \left(r|u'|^2+ \frac{k^2}{r}u^2 \right)
        + \frac{\Pe}2 \,k\, \xi' u^2 \, \dr \geq 0
\end{equation*}
for all $k \in \Z$ and all $u\in H^1(\rho_0,1)$ satisfying $u(1)=0$. One can verify that this inequality holds for our choice of $\xi$ and $\alpha$ with the same arguments used to prove \cref{prop:annulus-feasibility}, in particular using inequality \cref{e:annulusc-cnstr-2}. This concludes the proof of \cref{th:refined-disk}.
\section{Conclusion}\label{s:conclusion}

In this work, we used Lagrange duality and semidefinite programming to bound the maximum cooling efficiency of steady, energy-constrained incompressible flows in two dimensions. We measured efficiency via the inverse of the mean square gradient of the fluid's temperature $T$, so  optimal flows solve an optimal control problem minimizing $\smash{\int_\Omega \abs{\nabla T}^2 \dx}$. We showed that the Lagrange dual of this optimal control problem is a well-posed `infinite-dimensional SDP' that can be approximated with guaranteed convergence by finite-dimensional SDP discretizations. We demonstrated this on optimal cooling problems in a square and an annulus, showing also how problem symmetries can be exploited to pose SDPs with reduced complexity. Finally, we proved new analytical bounds on the optimal cooling efficiency $\efficiency(\Pe)$, where the square of the \Peclet\ number $\Pe$ measures the flow's kinetic energy. 
Specifically, we first extended to general domains the bound $\efficiency(\Pe)\lesssim\Pe^{2}$ proved for disks in \cite{Tobasco2022}. Then, inspired by computational results obtained with our semidefinite programming framework, we established the improved upper bound $\efficiency(\Pe) \lesssim \Pe^{2} / \ln^2\Pe$ for flows in disks and annuli when the heat distribution has positive azimuthal average.

Although we focused only on steady, energy-constrained flows in heated domains with a cold boundary, our techniques can be applied more generally. Examples include optimal cooling problems with a mix of isothermal and insulating boundaries, with no-slip boundary conditions, and with bounds on the $H^s$ norm of the fluid's velocity for $s\geq 1$. (The case $s=1$ corresponds to the enstrophy-constrained flows studied in \cite{Tobasco2022, Song2023, Iyer2022, DoeringTobasco2019}.) Problems in which the cooling efficiency $\int_\Omega |\nabla T|^2\,\dx$ is replaced by a generic bounded quadratic functional $\mathcal{E}(T)$ on $\Hoz$ can also be handled with little additional effort. Indeed, since solutions to the advection-diffusion equation \cref{e:ade} satisfy the `energy balance' $\int_\Omega |\nabla T|^2 + fT \,\dx = 0$, it suffices to bound the `translated' cooling efficiency $\mathcal{E}(T) + b\int_\Omega |\nabla T|^2 + fT \,\dx$ for $b>0$ large enough to make this functional bounded below (see also \cite[\S2.1]{Tobasco2022} for a similar approach). One can even view $b$ as a Lagrange multiplier for the energy balance and optimize it computationally within our semidefinite programming framework. Finally, our methods and the extensions we just described can be generalized to time-dependent optimal cooling problems where the cooling efficiency is defined in a space-time averaged sense. The main differences are additional technicalities in the definition of weak solutions to the advection-diffusion equation, but these can be handled as explained in \cite{Tobasco2022}.

Our work leaves open a number of questions. First, although we have proved that SDP discretizations of the dual problem \cref{e:ldp} converge, we have no convergence rates. Our computations suggest that classical approximation rates for Galerkin methods should apply. Proving this, however, requires a higher than $H^1$ regularity for the optimal dual variable $\xi$ and for the `critical fields' associated with the constraint of \cref{e:ldp}. One way to establish this higher regularity would be to study the optimality conditions for \cref{e:ldp}. 

A second open challenge is to remove the computational bottlenecks in the solution of high-resolution SDP discretizations of \cref{e:ldp}.
One option is to exploit structural properties of \cref{e:ldp} and its discretizations beyond symmetries. For example, the small number of `critical fields' observed in our computations suggests that the dual of the SDP \cref{e:sdp} generically has a low-rank solution. SDP solvers exploiting this low-rank structure (see, e.g., \cite{Burer2003,Burer2006,Bellavia2021,Habibi2023,Souto2022,Monteiro2026}) may thus be able to accurately approximate $\lb(\Pe)$ over a large enough range of \Peclet\ numbers to estimate the large-$\Pe$ asymptotics. Another option is to replace general-purpose algorithms for finite-dimensional SDP with function-space algorithms tailored to the infinite-dimensional problem \cref{e:ldp}. Unfortunately, we are not aware of any such algorithms beyond the gradient methods developed in \cite{Wen2013,Wen2015,Ding2019,Wen2022} for problems arising when bounding the time-averaged dissipation of turbulent flows, which can be adapted to \cref{e:ldp} but do not have general convergence guarantees. 

Finally, since the optimal cooling problem \cref{e:ocp} is not convex, the bound $\lb(\Pe)$ obtained with its Lagrange dual \cref{e:ldp} should not be sharp in general. While non-sharp bounds may still exhibit the same large-$\Pe$ asymptotics as the true optimal cost $\ocpval(\Pe)$, we do not know any nontrivial case in which this happens. (Of course, $\ocpval(\Pe)=\lb(\Pe)=0$ for all $\Pe$ if the domain is not heated.) If duality gaps affect the large-$\Pe$ asymptotics, or if quantitatively sharp bounds are required, then one must go beyond Lagrange duality. For discretizations of \cref{e:ocp}, which are quadratic programs, the moment-sum-of-squares hierarchy \cite{Laurent2009,Lasserre2010,Lasserre2024} provides a sequence of convex dual problems that generalize the Lagrange dual and provably converge if certain `algebraic compactness' conditions hold. In principle, it should be possible to formulate an infinite-dimensional version of this hierarchy by adapting ideas proposed in \cite{Goulart2012,Chernyshenko2014a,Goluskin2019,HenrionRudi2025,henrion2023infmom}. While current computational bottlenecks are likely to make a naive application of this approach prohibitively expensive, the particular structure of problem \cref{e:ocp} may allow for some progress. How much can be achieved remains to be seen, but it is an avenue worth exploring.

% Appendix
\appendix
\crefalias{section}{appendix}
\section{Fourier expansion of \texorpdfstring{\cref{e:ldp}}{(\ref{e:lpd})} in an annulus} \label{app:fourier-annulus}
%
Consider problem \cref{e:ldp} in the annulus $\Omega = \smash{\{\vec{x}:\;\rho < \abs{\vec{x}}^2 < 1\}}$. We work in polar coordinates $(r,\phi)$ and assume that $f=f(r)$ depends only on the radial direction. We can then take $\xi=\xi(r)$, too, so the objective function in \cref{e:ldp} reduces to that in~\cref{e:ldp-annulus-symm}. 

We claim that the constraints of these two problems are also equivalent. To show this, we introduce the Fourier series expansions
\begin{equation*}
    \psi(r,\phi) = \sum_{k\in\Z} \hat{\psi}_k(r) e^{ik\phi}
    \quad \text{and} \quad
    \theta(r,\phi) = \sum_{k\in\Z} \hat{\theta}_k(r) e^{ik\phi},
\end{equation*}
where the complex-valued functions $\smash{\hat{\psi}_k}$ and $\smash{\hat{\theta}_k}$ satisfy $\smash{\hat{\psi}_{-k}=\hat{\psi}_k^*}$ and $\smash{\hat{\theta}_{-k}=\hat{\theta}_k^*}$ because $\psi$ and $\theta$ are real-valued. The boundary conditions on $\psi$ and $\theta$ require $\smash{\hat{\theta}_0(\rho)}=0$, $\smash{\hat{\psi}_0(1)} = 0$, $\smash{\hat{\theta}_0(1)}=0$ and
$\smash{\hat{\psi}_k(\rho) = 0}$, $\smash{\hat{\theta}_k(\rho)}=0$, $\smash{\hat{\psi}_k(1)} = 0$, and $\smash{\hat{\theta}_k(1)}=0$ for all $k \neq 0$.
In particular, $\smash{\hat{\psi}_k,\hat{\theta}_k \in H^1_0(\rho,1)}$ for all $k \neq 0$.

Next, we compute the Fourier expansion of the quadratic form $\QF(a,\xi)[\psi,\theta]$ from \cref{e:qf}. Using the orthogonality of the Fourier modes, we find that
\begin{equation*}
    \frac{1}{2\pi}\int_\Omega |\nabla \psi|^2 \dx =
    \sum_{k\in \Z} \int_\rho^1 \left( |\hat{\psi}_k'|^2  + \frac{k^2}{r^2} |\hat{\psi}_k|^2 \right) r\dr.
\end{equation*}
A similar identity holds for $\int |\nabla \theta|^2 \dx$. For the term $\int \xi \langle \nabla^\perp \psi, \nabla \theta\rangle\,\dx$, instead, we first integrate by parts and then expand in Fourier modes to find that
\begin{equation*}
    \frac{1}{2\pi}\int \xi \langle \nabla^\perp \psi, \nabla \theta\rangle\,\dx
    = 
    -\frac{1}{2\pi}\int \theta \langle \nabla^\perp \psi, \nabla \xi\rangle\,\dx
    = 
    \sum_{k\in \Z}  \int_\rho^1  ik \xi' {\hat{\psi}}_k^* \hat{\theta}_{k} \xi' \dr.
\end{equation*}
Upon combining everything, we obtain
\begin{equation*}
    \frac{1}{2\pi}\QF(a,\xi)[\psi,\theta]
    =\sum_{k \in \Z} 
    \int_\rho^1
    a \left( r|\hat{\psi}_k'|^2  + \frac{k^2}{r} |\hat{\psi}_k|^2 \right)
    + \left( r|\hat{\theta}_k'|^2  + \frac{k^2}{r} |\hat{\theta}_k|^2 \right)
    + ik \xi' {\hat{\psi}}_k^* \hat{\theta}_{k} \xi' \dr.
\end{equation*}

Now, write $\psi_k = A_k - iB_k$ and $\theta_k = C_k+iD_k$ for real-valued functions $A_k$, $B_k$, $C_k$, and $D_k$ in $H^1_0(\rho,1)$. Let $\QF_k(a,\xi)$ be the quadratic form in \cref{e:qf-fourier-annulus}. A straightforward calculation yields
\begin{align*}
    \frac{1}{2\pi}\QF(a,\xi)[\psi,\theta]
    &=\sum_{k \in \Z} 
    \QF_k(a,\xi)[A_k,D_k] +
    \QF_k(a,\xi)[B_k,C_k].
\end{align*}
Since $A_k,B_k,C_k,D_k \in H^1_0(\rho,1)$ can be chosen independently, the condition $\QF(a,\xi)\succeq 0$ holds on $\cal{H}\times\Hoz$ if and only if $\QF_k(a,\xi) \succeq 0$ on $H^1_0(\rho,1) \times H^1_0(\rho,1)$ for all $k\in \Z$. This proves our claim that the constraint of \cref{e:ldp} and \cref{e:ldp-annulus-symm} are equivalent when $\Omega$ is an annulus and the heat distribution $f$ is radially symmetric.

% % Acknowledgments
% \subsection*{Acknowledgments}
% Write acknowledgments here if necessary.

% Bibliography
\bibliographystyle{abbrvnat}
\bibliography{reflist}

%%%%%%%%%%%%%%%%%%%%%%%%%%%%%%%%%%%%%%%%%%%%%%%%%%%%%%%%%%%%%%%%%%%%%%%%%%%%%%%%
%%%%%%%%%%%%%%%%%%%%%%%%%%%%%%%%%%%%%%%%%%%%%%%%%%%%%%%%%%%%%%%%%%%%%%%%%%%%%%%%
\end{document}
